% =====================================================================
\chapter{Soluciones de los ejercicios}
\label{ap:soluciones}
% =====================================================================

Cada solución se presenta en cinco partes: datos, planteamiento, desarrollo, resultado y comprobación. En los ejercicios de demostración, los datos son las hipótesis y el resultado es el enunciado demostrado; la comprobación consiste en un caso particular o en una verificación independiente. Los recuentos se han comprobado por enumeración exhaustiva con un programa.

% ---------------------------------------------------------------------
\section{Tema 1}

\begin{solucion}{eje:1-pert}
\datos $A=\{1,\{1\},\{1,2\}\}$, cuyos tres elementos son $1$, $\{1\}$ y $\{1,2\}$.
\planteamiento $x\in A$ exige que $x$ sea uno de los tres elementos; $X\subseteq A$ exige que \emph{cada} elemento de $X$ sea uno de ellos.
\pasos
\begin{enumerate}[label=(\alph*)]
  \item $1$ es elemento de $A$: \textbf{verdadera}.
  \item $\{1\}$ es elemento de $A$: \textbf{verdadera}.
  \item El único elemento de $\{1\}$ es $1\in A$: \textbf{verdadera}.
  \item El elemento $2$ de $\{2\}$ no está en $A$: \textbf{falsa}.
  \item $\{1,2\}$ es elemento de $A$: \textbf{verdadera}.
  \item El único elemento de $\{\{1,2\}\}$ es $\{1,2\}\in A$: \textbf{verdadera}.
  \item $\emptyset$ es subconjunto de todo conjunto (\cref{prop:t1-subconj}): \textbf{verdadera}.
  \item $\emptyset$ no es ninguno de los tres elementos: \textbf{falsa}.
\end{enumerate}
\resultado verdaderas (a), (b), (c), (e), (f), (g); falsas (d), (h).
\comprobacion los apartados (b) y (c) muestran que un mismo objeto, $\{1\}$, puede ser a la vez elemento y subconjunto de $A$.
\end{solucion}

\begin{solucion}{eje:1-card}
\planteamiento contar los elementos distintos, sin mirar dentro de los elementos que son conjuntos.
\pasos
\begin{enumerate}[label=(\alph*)]
  \item $\{\emptyset\}$ tiene un elemento: $1$.
  \item $\emptyset$ y $\{\emptyset\}$ son distintos (uno vacío, otro no): $2$.
  \item $\Pa(\emptyset)=\{\emptyset\}$: $1=2^{0}$.
  \item $\Pa(\Pa(\emptyset))=\Pa(\{\emptyset\})=\{\emptyset,\{\emptyset\}\}$: $2=2^{1}$.
  \item $\{1,2\}=\{2,1\}$, luego el conjunto es $\{\{1,2\}\}$: $1$.
  \item $1$, $\{1\}$ y $\{\{1\}\}$ son distintos: $3$.
\end{enumerate}
\resultado $1,\ 2,\ 1,\ 2,\ 1,\ 3$.
\comprobacion (c) y (d) concuerdan con el \cref{teo:t1-partes}: $\abs{\Pa(A)}=2^{\abs{A}}$ con $\abs{A}=0$ y $\abs{A}=1$.
\end{solucion}

\begin{solucion}{eje:1-extension}
\datos $A=\{x\in\Z: x^{2}<10\}$, $B$ = divisores positivos de 12, $C$ = raíces de $x^{2}-3x+2$.
\planteamiento escribir cada conjunto por extensión y aplicar las definiciones.
\pasos
\begin{enumerate}[label=\arabic*.]
  \item $x^{2}<10$ con $x$ entero: $\abs{x}\le3$ porque $3^{2}=9<10\le16=4^{2}$. $A=\{-3,-2,-1,0,1,2,3\}$.
  \item $B=\{1,2,3,4,6,12\}$.
  \item $x^{2}-3x+2=(x-1)(x-2)$: $C=\{1,2\}$.
  \item $A\cap B=\{1,2,3\}$; $A\setminus B=\{-3,-2,-1,0\}$; $B\setminus A=\{4,6,12\}$.
  \item $A\dsim B=(A\setminus B)\cup(B\setminus A)=\{-3,-2,-1,0,4,6,12\}$.
  \item $(A\cap B)\setminus C=\{1,2,3\}\setminus\{1,2\}=\{3\}$.
\end{enumerate}
\resultado $A\cap B=\{1,2,3\}$, $A\setminus B=\{-3,-2,-1,0\}$, $B\setminus A=\{4,6,12\}$, $A\dsim B=\{-3,-2,-1,0,4,6,12\}$, $(A\cap B)\setminus C=\{3\}$.
\comprobacion $\abs{A\dsim B}=\abs{A}+\abs{B}-2\abs{A\cap B}=7+6-6=7$.
\end{solucion}

\begin{solucion}{eje:1-partes}
\datos $A=\{a,b,c\}$, $\abs{A}=3$.
\planteamiento listar por tamaños y contar.
\pasos
\begin{enumerate}[label=\arabic*.]
  \item $\Pa(A)=\{\emptyset,\{a\},\{b\},\{c\},\{a,b\},\{a,c\},\{b,c\},\{a,b,c\}\}$: 8 elementos.
  \item Contienen a $a$: $\{a\},\{a,b\},\{a,c\},\{a,b,c\}$, es decir, 4.
  \item Con dos elementos: $\{a,b\},\{a,c\},\{b,c\}$, es decir, 3.
  \item En general, $X\mapsto X\setminus\{a\}$ hace corresponder los subconjuntos que contienen a $a$ con los subconjuntos de $A\setminus\{a\}$, uno a uno (la operación inversa es añadir $a$). Hay $2^{n-1}$.
\end{enumerate}
\resultado 8 subconjuntos; 4 contienen a $a$; 3 tienen dos elementos; en general, $2^{n-1}$ contienen a $a$.
\comprobacion $2^{3-1}=4$; además los que no contienen a $a$ son otros $4$ y $4+4=8$.
\end{solucion}

\begin{solucion}{eje:1-cartesiano}
\datos $A=\{0,1\}$, $B=\{\cad{a},\cad{b},\cad{c}\}$.
\planteamiento \cref{prop:t1-producto}.
\pasos
\begin{enumerate}[label=(\alph*)]
  \item $A\times B=\{(0,\cad{a}),(0,\cad{b}),(0,\cad{c}),(1,\cad{a}),(1,\cad{b}),(1,\cad{c})\}$.
  \item $\abs{A^{3}}=2^{3}=8$; $\abs{B\times A\times A}=3\cdot2\cdot2=12$.
  \item No: $(0,\cad{a})\in A\times B$, pero $(0,\cad{a})\notin B\times A$, cuyos pares tienen primera componente en $B$.
  \item $A^{8}$ es el conjunto de las cadenas de 8 bits, es decir, de los bytes: $\abs{A^{8}}=2^{8}=256$.
\end{enumerate}
\resultado 6 pares; $8$ y $12$; $A\times B\neq B\times A$; hay $256$ bytes distintos.
\comprobacion $\abs{A\times B}=2\cdot3=6$ coincide con la lista de~(a).
\end{solucion}

\begin{solucion}{eje:1-demorgan-num}
\datos $U=\{1,\dots,10\}$, $A=\{1,2,3,4,5\}$, $B=\{2,4,6,8,10\}$, $C=\{3,4,5,6\}$.
\planteamiento calcular cada miembro por separado.
\pasos
\begin{enumerate}[label=\arabic*.]
  \item $A\cup B=\{1,2,3,4,5,6,8,10\}$, luego $\overline{A\cup B}=\{7,9\}$.
  \item $\overline{A}=\{6,7,8,9,10\}$, $\overline{B}=\{1,3,5,7,9\}$, luego $\overline{A}\cap\overline{B}=\{7,9\}$.
  \item $B\cap C=\{4,6\}$, luego $A\setminus(B\cap C)=\{1,2,3,5\}$.
  \item $A\setminus B=\{1,3,5\}$ y $A\setminus C=\{1,2\}$, luego $(A\setminus B)\cup(A\setminus C)=\{1,2,3,5\}$.
  \item Coinciden porque, por el \cref{teo:t1-leyes}, $A\setminus(B\cap C)=A\cap\overline{B\cap C}=A\cap(\overline{B}\cup\overline{C})=(A\cap\overline{B})\cup(A\cap\overline{C})$, que es $(A\setminus B)\cup(A\setminus C)$.
\end{enumerate}
\resultado $\overline{A\cup B}=\overline{A}\cap\overline{B}=\{7,9\}$; $A\setminus(B\cap C)=(A\setminus B)\cup(A\setminus C)=\{1,2,3,5\}$.
\comprobacion $7$ y $9$ son los únicos números de $U$ que no son pares ni están entre 1 y~5.
\end{solucion}

\begin{solucion}{eje:1-vf}
\planteamiento buscar contraejemplos pequeños; si no aparecen, intentar una demostración.
\pasos
\begin{enumerate}[label=(\alph*)]
  \item \textbf{Falsa.} Con $A=\{1,2\}$, $B=\{2\}$: $(A\cup B)\setminus B=\{1\}\neq A$.
  \item \textbf{Verdadera.} Es el \cref{ej:t1-doble}.
  \item \textbf{Falsa.} Con $A=B=\{1\}$ y $C=D=\{2\}$: el miembro izquierdo es $\{(1,1),(2,2)\}$; el derecho, $\{1,2\}\times\{1,2\}$, contiene además $(1,2)$ y $(2,1)$. Solo se cumple la inclusión $\subseteq$.
\end{enumerate}
\resultado (a) falsa; (b) verdadera; (c) falsa.
\comprobacion en (c), los cardinales $2$ y $4$ ya muestran que los conjuntos son distintos.
\end{solucion}

\begin{solucion}{eje:1-dif-union}
\datos $A,B,C$ conjuntos arbitrarios.
\planteamiento doble inclusión (\cref{prop:t1-subconj}(iii)).
\pasos
\begin{enumerate}[label=\arabic*.]
  \item ($\subseteq$) Sea $x\in A\setminus(B\cup C)$. Entonces $x\in A$ y $x\notin B\cup C$, es decir, $x\notin B$ y $x\notin C$. Por tanto $x\in A\setminus B$ y $x\in A\setminus C$, y $x$ está en la intersección.
  \item ($\supseteq$) Sea $x\in(A\setminus B)\cap(A\setminus C)$. Entonces $x\in A$, $x\notin B$ y $x\notin C$; luego $x\notin B\cup C$ y $x\in A\setminus(B\cup C)$.
\end{enumerate}
\resultado $A\setminus(B\cup C)=(A\setminus B)\cap(A\setminus C)$. \hfill$\blacksquare$
\comprobacion con los datos del \cref{eje:1-demorgan-num}: $A\setminus(B\cup C)=\{1\}$ y $(A\setminus B)\cap(A\setminus C)=\{1,3,5\}\cap\{1,2\}=\{1\}$.
\end{solucion}

\begin{solucion}{eje:1-union-inclusion}
\datos $A,B$ conjuntos arbitrarios.
\planteamiento probar las dos implicaciones; la igualdad $A\cup B=B$ se prueba por doble inclusión.
\pasos
\begin{enumerate}[label=\arabic*.]
  \item ($\Rightarrow$) Supóngase $A\subseteq B$. Siempre $B\subseteq A\cup B$. Recíprocamente, si $x\in A\cup B$, entonces $x\in A$ o $x\in B$; en el primer caso $x\in B$ por la hipótesis. Luego $A\cup B\subseteq B$ y hay igualdad.
  \item ($\Leftarrow$) Supóngase $A\cup B=B$. Si $x\in A$, entonces $x\in A\cup B=B$. Luego $A\subseteq B$.
  \item Análogo para la intersección: $A\subseteq B\iff A\cap B=A$. ($\Rightarrow$) $A\cap B\subseteq A$ siempre; si $x\in A$, entonces $x\in B$ por hipótesis y $x\in A\cap B$. ($\Leftarrow$) Si $x\in A=A\cap B$, entonces $x\in B$.
\end{enumerate}
\resultado $A\subseteq B\iff A\cup B=B\iff A\cap B=A$. \hfill$\blacksquare$
\comprobacion con $A=\{1\}$, $B=\{1,2\}$: $A\cup B=\{1,2\}=B$ y $A\cap B=\{1\}=A$.
\end{solucion}

\begin{solucion}{eje:1-producto-inter}
\datos $A,B,C$ conjuntos arbitrarios.
\planteamiento los elementos de un producto son pares; se trabaja con un par arbitrario $(x,y)$.
\pasos
\begin{enumerate}[label=\arabic*.]
  \item $(x,y)\in A\times(B\cap C)\iff x\in A$ y $y\in B$ y $y\in C\iff(x\in A$ y $y\in B)$ y $(x\in A$ y $y\in C)\iff(x,y)\in(A\times B)\cap(A\times C)$. Cada paso es una equivalencia, luego se tienen las dos inclusiones a la vez.
  \item ($\subseteq$) Sea $(x,y)\in(A\setminus B)\times C$: $x\in A$, $x\notin B$, $y\in C$. Entonces $(x,y)\in A\times C$, y $(x,y)\notin B\times C$ porque $x\notin B$.
  \item ($\supseteq$) Sea $(x,y)\in(A\times C)\setminus(B\times C)$: $x\in A$, $y\in C$ y no es cierto que ($x\in B$ y $y\in C$). Como $y\in C$, debe ser $x\notin B$. Luego $(x,y)\in(A\setminus B)\times C$.
\end{enumerate}
\resultado $A\times(B\cap C)=(A\times B)\cap(A\times C)$ y $(A\setminus B)\times C=(A\times C)\setminus(B\times C)$. \hfill$\blacksquare$
\comprobacion con $A=\{1,2\}$, $B=\{2\}$, $C=\{0\}$: $(A\setminus B)\times C=\{(1,0)\}$ y $(A\times C)\setminus(B\times C)=\{(1,0),(2,0)\}\setminus\{(2,0)\}=\{(1,0)\}$.
\end{solucion}

\begin{solucion}{eje:1-simetrica}
\datos $A,B,C$ conjuntos; $A\dsim B=(A\setminus B)\cup(B\setminus A)$.
\planteamiento (a) doble inclusión; (b) lectura de la definición; (c) traducir la pertenencia a paridades con funciones características (\cref{tema:2}).
\pasos
\begin{enumerate}[label=\arabic*.]
  \item (a, $\subseteq$) Si $x\in A\setminus B$, entonces $x\in A\cup B$ y $x\notin A\cap B$ (no está en $B$); igual si $x\in B\setminus A$.
  \item (a, $\supseteq$) Si $x\in A\cup B$ y $x\notin A\cap B$: si $x\in A$, entonces $x\notin B$ y $x\in A\setminus B$; si $x\in B$, entonces $x\notin A$ y $x\in B\setminus A$.
  \item (b) $x\in A\setminus B$ significa «en $A$ y no en $B$», y $x\in B\setminus A$, «en $B$ y no en $A$»: $x\in A\dsim B$ exactamente cuando está en uno y solo uno.
  \item (c) Por (b), $\chi_{A\dsim B}(x)=1$ si y solo si $\chi_A(x)+\chi_B(x)=1$, es decir, $\chi_{A\dsim B}=(\chi_A+\chi_B)\bmod 2$.
  \item Aplicándolo dos veces: $\chi_{(A\dsim B)\dsim C}=\bigl((\chi_A+\chi_B)\bmod2+\chi_C\bigr)\bmod2=(\chi_A+\chi_B+\chi_C)\bmod 2$. Del mismo modo, $\chi_{A\dsim(B\dsim C)}=(\chi_A+\chi_B+\chi_C)\bmod2$.
  \item Los dos conjuntos tienen la misma función característica, luego son iguales: ambos son los $x$ que pertenecen a un número \emph{impar} de los tres conjuntos.
\end{enumerate}
\resultado $A\dsim B=(A\cup B)\setminus(A\cap B)$ y $(A\dsim B)\dsim C=A\dsim(B\dsim C)$. \hfill$\blacksquare$
\comprobacion con $A=\{1,2\}$, $B=\{2,3\}$, $C=\{3,4\}$: $(A\dsim B)\dsim C=\{1,3\}\dsim\{3,4\}=\{1,4\}$ y $A\dsim(B\dsim C)=\{1,2\}\dsim\{2,4\}=\{1,4\}$.
\end{solucion}

\begin{solucion}{eje:1-partes-inter}
\datos $A,B$ conjuntos arbitrarios.
\planteamiento los elementos de $\Pa(\cdot)$ son conjuntos: se toma un conjunto arbitrario $X$.
\pasos
\begin{enumerate}[label=\arabic*.]
  \item $X\in\Pa(A)\cap\Pa(B)\iff X\subseteq A$ y $X\subseteq B$.
  \item Si $X\subseteq A$ y $X\subseteq B$, todo $x\in X$ está en $A$ y en $B$, luego $X\subseteq A\cap B$. Recíprocamente, si $X\subseteq A\cap B$, como $A\cap B\subseteq A$ y $A\cap B\subseteq B$, la transitividad da $X\subseteq A$ y $X\subseteq B$. Así $X\in\Pa(A)\cap\Pa(B)\iff X\in\Pa(A\cap B)$.
  \item Si $X\in\Pa(A)\cup\Pa(B)$, entonces $X\subseteq A$ o $X\subseteq B$; en ambos casos $X\subseteq A\cup B$, luego $X\in\Pa(A\cup B)$.
  \item Contraejemplo de la igualdad: $A=\{1\}$, $B=\{2\}$. Entonces $\{1,2\}\in\Pa(A\cup B)$, pero $\{1,2\}\not\subseteq A$ y $\{1,2\}\not\subseteq B$.
\end{enumerate}
\resultado $\Pa(A)\cap\Pa(B)=\Pa(A\cap B)$ y $\Pa(A)\cup\Pa(B)\subseteq\Pa(A\cup B)$, con inclusión estricta en general. \hfill$\blacksquare$
\comprobacion en el contraejemplo, $\abs{\Pa(A)\cup\Pa(B)}=\abs{\{\emptyset,\{1\},\{2\}\}}=3<4=\abs{\Pa(A\cup B)}$.
\end{solucion}

\begin{solucion}{eje:1-familias}
\datos $B_n=\{k\in\N: k\ge n\}$; intervalos $(\tfrac1n,2)$ y $(-\tfrac1n,\tfrac1n)$.
\planteamiento para la intersección infinita, suponer que un elemento está en todos los conjuntos y llegar a contradicción.
\pasos
\begin{enumerate}[label=\arabic*.]
  \item Los $B_n$ son decrecientes ($B_{n+1}\subseteq B_n$), luego $B_1\cap\dots\cap B_n=B_n$, que contiene a $n$: no es vacío.
  \item Si $k\in\bigcap_{n\ge1}B_n$, en particular $k\in B_{k+1}$, es decir, $k\ge k+1$: absurdo. Luego $\bigcap_{n\ge1}B_n=\emptyset$.
  \item Cada $(\tfrac1n,2)\subseteq(0,2)$. Si $x\in(0,2)$, se elige $n>1/x$; entonces $\tfrac1n<x$ y $x\in(\tfrac1n,2)$. Luego $\bigcup_{n\ge1}(\tfrac1n,2)=(0,2)$.
  \item $0$ está en todos los $(-\tfrac1n,\tfrac1n)$. Si $x\neq0$, se elige $n>1/\abs{x}$; entonces $\abs{x}>\tfrac1n$ y $x$ no está en el $n$-ésimo intervalo. Luego la intersección es $\{0\}$.
\end{enumerate}
\resultado $\bigcap B_n=\emptyset$ aunque toda intersección finita es no vacía; $\bigcup(\tfrac1n,2)=(0,2)$; $\bigcap(-\tfrac1n,\tfrac1n)=\{0\}$. \hfill$\blacksquare$
\comprobacion $0\notin\bigcup(\tfrac1n,2)$ porque $0<\tfrac1n$ para todo $n$, de acuerdo con el extremo abierto de $(0,2)$.
\end{solucion}

\begin{solucion}{eje:1-cadenas}
\datos $\Sigma=\{\cad{a},\cad{b}\}$; $L_1$: longitud $\le2$; $L_2$: empiezan por \cad{a}.
\planteamiento listar por longitudes.
\pasos
\begin{enumerate}[label=(\alph*)]
  \item $L_1=\{\vacia,\cad{a},\cad{b},\cad{aa},\cad{ab},\cad{ba},\cad{bb}\}$; $\abs{L_1}=1+2+4=7$.
  \item $L_1\cap L_2=\{\cad{a},\cad{aa},\cad{ab}\}$; $L_1\setminus L_2=\{\vacia,\cad{b},\cad{ba},\cad{bb}\}$ (la cadena vacía no empieza por \cad{a}).
  \item $L_2$ es infinito: contiene $\cad{a},\cad{aa},\cad{aaa},\dots$, todas distintas.
\end{enumerate}
\resultado $\abs{L_1}=7$; $L_1\cap L_2$ tiene 3 cadenas y $L_1\setminus L_2$ tiene~4; $L_2$ es infinito.
\comprobacion $3+4=7=\abs{L_1}$, como debe ser al partir $L_1$ en dos.
\end{solucion}

% ---------------------------------------------------------------------
\section{Tema 2}

\begin{solucion}{eje:2-esfuncion}
\planteamiento comprobar existencia y unicidad de la imagen (\cref{def:t2-funcion}) y, en (e), que la regla no dependa del representante.
\pasos
\begin{enumerate}[label=(\alph*)]
  \item Cada uno de $1,2,3$ aparece exactamente una vez: \textbf{es función}.
  \item El~3 no tiene imagen: \textbf{no es función} (falla la existencia).
  \item El~1 tiene dos imágenes: \textbf{no es función} (falla la unicidad).
  \item $0$ no tiene imagen en $\R$: \textbf{no es función} de $\R$ en $\R$ (sí lo es de $\R\setminus\{0\}$ en $\R$).
  \item $\tfrac12=\tfrac24$, pero $\tfrac{1+1}{2+1}=\tfrac23$ y $\tfrac{2+1}{4+1}=\tfrac35$: \textbf{no está bien definida}.
\end{enumerate}
\resultado solo (a) es función.
\comprobacion en (e), $\tfrac23\neq\tfrac35$ porque $2\cdot5=10\neq9=3\cdot3$.
\end{solucion}

\begin{solucion}{eje:2-tabla}
\datos $f(0)=2$, $f(1)=4$, $f(2)=2$, $f(3)=1$, $f(4)=2$; $B=\{1,2,3,4,5\}$.
\planteamiento \cref{def:t2-imagen}.
\pasos
\begin{enumerate}[label=\arabic*.]
  \item Cada elemento de $A$ aparece en un único par y las segundas componentes están en $B$: es función. Dominio $A$, codominio $B$, recorrido $f(A)=\{1,2,4\}$.
  \item $f(\{0,1,2\})=\{f(0),f(1),f(2)\}=\{2,4\}$.
  \item $f^{-1}(\{2\})=\{0,2,4\}$; $f^{-1}(\{3,5\})=\emptyset$; $f^{-1}(\{1,4\})=\{1,3\}$.
\end{enumerate}
\resultado recorrido $\{1,2,4\}$; $f(\{0,1,2\})=\{2,4\}$; preimágenes $\{0,2,4\}$, $\emptyset$, $\{1,3\}$.
\comprobacion las fibras $f^{-1}(\{1\})=\{3\}$, $f^{-1}(\{2\})=\{0,2,4\}$, $f^{-1}(\{4\})=\{1\}$ tienen $1+3+1=5=\abs{A}$ elementos.
\end{solucion}

\begin{solucion}{eje:2-contar}
\planteamiento \cref{prop:t2-contar}: $\abs{B^{A}}=\abs{B}^{\abs{A}}$.
\pasos
\begin{enumerate}[label=\arabic*.]
  \item De 3 en 4: $4^{3}=64$. De 4 en 3: $3^{4}=81$.
  \item $\abs{\{0,1\}^{2}}=4$, luego hay $2^{4}=16$ funciones booleanas de dos variables.
  \item $\abs{\{0,1\}^{n}}=2^{n}$, luego hay $2^{2^{n}}$; para $n=3$, $2^{8}=256$.
\end{enumerate}
\resultado $64$, $81$, $16$ y $2^{2^{n}}$ ($256$ para $n=3$).
\comprobacion entre las 16 funciones booleanas están las conocidas \textsc{and}, \textsc{or} y \textsc{xor}, y las dos constantes.
\end{solucion}

\begin{solucion}{eje:2-cuadrado}
\datos $f(x)=x^{2}-2x=(x-1)^{2}-1$.
\planteamiento completar el cuadrado: $f$ decrece en $(-\infty,1]$, crece en $[1,\infty)$ y tiene mínimo $-1$ en $x=1$.
\pasos
\begin{enumerate}[label=\arabic*.]
  \item En $[0,3]$: mínimo $f(1)=-1$; $f(0)=0$ y $f(3)=3$. Por continuidad, $f([0,3])=[-1,3]$.
  \item $f(x)=3\iff x^{2}-2x-3=0\iff(x-3)(x+1)=0$: $f^{-1}(\{3\})=\{-1,3\}$.
  \item $-1\le(x-1)^{2}-1\le0\iff0\le(x-1)^{2}\le1\iff\abs{x-1}\le1$: $f^{-1}([-1,0])=[0,2]$.
  \item $f(x)\ge-1>-2$ para todo $x$: $f^{-1}((-\infty,-2))=\emptyset$.
  \item $f(\{3\})=\{3\}$, luego $f^{-1}(f(\{3\}))=\{-1,3\}\supsetneq\{3\}$.
\end{enumerate}
\resultado $[-1,3]$; $\{-1,3\}$; $[0,2]$; $\emptyset$; $\{-1,3\}$.
\comprobacion $f(-1)=1+2=3$, y $f(0)=f(2)=0$ son los extremos de~3.
\end{solucion}

\begin{solucion}{eje:2-suelo}
\datos $f(n)=\lfloor n/2\rfloor$ en $\Z$.
\planteamiento $f(2k)=f(2k+1)=k$ para todo $k\in\Z$.
\pasos
\begin{enumerate}[label=\arabic*.]
  \item $f(0)=0$, $f(1)=0$, $f(2)=1$, $f(3)=1$: $f(\{0,1,2,3\})=\{0,1\}$.
  \item $f^{-1}(\{1\})=\{2,3\}$, y $f^{-1}(\{0,1,2\})=\{0,1,2,3,4,5\}$.
  \item $f(2\Z)=\{f(2k):k\in\Z\}=\{k:k\in\Z\}=\Z$.
  \item Como ya $f(2\Z)=\Z$, el recorrido es $\Z$.
\end{enumerate}
\resultado $\{0,1\}$; $\{2,3\}$; $\{0,\dots,5\}$; $\Z$; recorrido $\Z$.
\comprobacion cada fibra $f^{-1}(\{k\})=\{2k,2k+1\}$ tiene dos elementos.
\end{solucion}

\begin{solucion}{eje:2-longitud}
\datos $\ell\colon\{\cad{a},\cad{b}\}^{*}\to\N$.
\planteamiento la fibra $\ell^{-1}(\{n\})$ es $\Sigma^{n}$.
\pasos
\begin{enumerate}[label=\arabic*.]
  \item $\ell^{-1}(\{2\})=\{\cad{aa},\cad{ab},\cad{ba},\cad{bb}\}$.
  \item $\ell^{-1}(\{0,1\})=\{\vacia,\cad{a},\cad{b}\}$.
  \item $\ell(\{\cad{ab},\cad{ba},\cad{aab}\})=\{2,3\}$.
  \item Para cada $n$, $\cad{a}^{n}$ tiene longitud $n$: el recorrido es $\N$.
  \item $\abs{\ell^{-1}(\{n\})}=\abs{\Sigma^{n}}=2^{n}$.
\end{enumerate}
\resultado las fibras tienen $2^{n}$ elementos y el recorrido es $\N$.
\comprobacion para $n=2$, la lista del paso~1 tiene $4=2^{2}$ cadenas.
\end{solucion}

\begin{solucion}{eje:2-recorrido}
\planteamiento resolver $f(x)=y$ y ver para qué $y$ hay solución.
\pasos
\begin{enumerate}[label=(\alph*)]
  \item $e^{x}=y$ tiene solución $x=\ln y$ si y solo si $y>0$: recorrido $(0,\infty)$.
  \item $\operatorname{sen}$ toma todos los valores de $[-1,1]$ en $[-\pi/2,\pi/2]$ y ninguno fuera: $[-1,1]$.
  \item $h(x)=(x+1)^{2}+2\ge2$, y $h(x)=y$ con $y\ge2$ tiene la solución $x=-1+\sqrt{y-2}$: $[2,\infty)$.
  \item $1+x^{2}\ge1$, luego $0<k(x)\le1$; dado $y\in(0,1]$, $x=\sqrt{1/y-1}$ cumple $k(x)=y$: $(0,1]$.
\end{enumerate}
\resultado $(0,\infty)$; $[-1,1]$; $[2,\infty)$; $(0,1]$.
\comprobacion $h(-1)=2$ y $k(0)=1$ alcanzan los extremos cerrados.
\end{solucion}

\begin{solucion}{eje:2-preim-ui}
\datos $f\colon X\to Y$; $C,D\subseteq Y$.
\planteamiento cadena de equivalencias para un $x\in X$ arbitrario.
\pasos
\begin{enumerate}[label=\arabic*.]
  \item $x\in f^{-1}(C\cap D)\iff f(x)\in C\cap D\iff f(x)\in C$ y $f(x)\in D\iff x\in f^{-1}(C)$ y $x\in f^{-1}(D)\iff x\in f^{-1}(C)\cap f^{-1}(D)$.
  \item Si $C\subseteq D$ y $x\in f^{-1}(C)$, entonces $f(x)\in C\subseteq D$, luego $x\in f^{-1}(D)$.
\end{enumerate}
\resultado $f^{-1}(C\cap D)=f^{-1}(C)\cap f^{-1}(D)$, y la preimagen conserva la inclusión. \hfill$\blacksquare$
\comprobacion con $f(x)=x^{2}$, $C=[0,1]$, $D=[1,4]$: $f^{-1}(C\cap D)=f^{-1}(\{1\})=\{-1,1\}$ y $[-1,1]\cap([-2,-1]\cup[1,2])=\{-1,1\}$.
\end{solucion}

\begin{solucion}{eje:2-imagen-ui}
\datos $f\colon X\to Y$; $A,B\subseteq X$.
\planteamiento doble inclusión para la unión; contraejemplo con $f(x)=x^{2}$ para la diferencia.
\pasos
\begin{enumerate}[label=\arabic*.]
  \item ($\subseteq$) Si $y\in f(A\cup B)$, existe $x\in A\cup B$ con $f(x)=y$; si $x\in A$, $y\in f(A)$, y si $x\in B$, $y\in f(B)$.
  \item ($\supseteq$) Si $y\in f(A)$, existe $x\in A\subseteq A\cup B$ con $f(x)=y$, luego $y\in f(A\cup B)$; igual si $y\in f(B)$.
  \item Contraejemplo: $f(x)=x^{2}$, $A=\{-1,1\}$, $B=\{1\}$. $f(A\setminus B)=f(\{-1\})=\{1\}$, pero $f(A)\setminus f(B)=\{1\}\setminus\{1\}=\emptyset$.
  \item Inclusión válida: $f(A)\setminus f(B)\subseteq f(A\setminus B)$. Si $y\in f(A)$ e $y\notin f(B)$, existe $x\in A$ con $f(x)=y$; si fuera $x\in B$, se tendría $y\in f(B)$. Luego $x\in A\setminus B$ e $y\in f(A\setminus B)$.
\end{enumerate}
\resultado $f(A\cup B)=f(A)\cup f(B)$; en general solo $f(A)\setminus f(B)\subseteq f(A\setminus B)$. \hfill$\blacksquare$
\comprobacion en el contraejemplo, $\emptyset\subseteq\{1\}$, de acuerdo con el paso~4.
\end{solucion}

\begin{solucion}{eje:2-ida-vuelta}
\datos $f\colon X\to Y$, $A\subseteq X$, $C\subseteq Y$.
\planteamiento aplicar las definiciones de imagen y preimagen.
\pasos
\begin{enumerate}[label=\arabic*.]
  \item Si $a\in A$, entonces $f(a)\in f(A)$, luego $a\in f^{-1}(f(A))$.
  \item Si $y\in f(f^{-1}(C))$, existe $x\in f^{-1}(C)$ con $y=f(x)$; como $f(x)\in C$, $y\in C$.
  \item Ejemplos estrictos con $f\colon\R\to\R$, $f(x)=x^{2}$: para $A=[0,1]$, $f^{-1}(f(A))=f^{-1}([0,1])=[-1,1]\supsetneq A$; para $C=[-1,1]$, $f(f^{-1}(C))=f([-1,1])=[0,1]\subsetneq C$.
\end{enumerate}
\resultado $A\subseteq f^{-1}(f(A))$ y $f(f^{-1}(C))\subseteq C$, con desigualdad estricta en general. \hfill$\blacksquare$
\comprobacion la primera es igualdad si $f$ es inyectiva y la segunda si $f$ es sobreyectiva; $x^{2}$ en $\R$ no es ninguna de las dos cosas.
\end{solucion}

\begin{solucion}{eje:2-preim-comp}
\datos $f\colon X\to Y$; $C\subseteq Y$.
\planteamiento cadena de equivalencias.
\pasos
\begin{enumerate}[label=\arabic*.]
  \item $x\in f^{-1}(Y\setminus C)\iff f(x)\in Y\setminus C\iff f(x)\notin C\iff x\notin f^{-1}(C)\iff x\in X\setminus f^{-1}(C)$ (se usa que $f(x)\in Y$ siempre y que $x\in X$).
  \item $x\in f^{-1}(\emptyset)\iff f(x)\in\emptyset$, que nunca ocurre: $f^{-1}(\emptyset)=\emptyset$.
  \item $f(x)\in Y$ para todo $x\in X$: $f^{-1}(Y)=X$.
\end{enumerate}
\resultado la preimagen conserva complemento, vacío y total. \hfill$\blacksquare$
\comprobacion junto con el \cref{teo:t2-preimagen} y el \cref{eje:2-preim-ui}, la preimagen conserva todas las operaciones de conjuntos.
\end{solucion}

\begin{solucion}{eje:2-caracteristica}
\datos $X,Z\subseteq U$; $\chi$ toma valores en $\{0,1\}$.
\planteamiento cuatro casos según $(\chi_X(u),\chi_Z(u))$.
\pasos
\begin{enumerate}[label=\arabic*.]
  \item Tabla: $(0,0)$: $u\notin X\cap Z$, $u\notin X\cup Z$; producto $0$, fórmula de la unión $0$. $(1,0)$ y $(0,1)$: $u\notin X\cap Z$, $u\in X\cup Z$; producto $0$, unión $1+0-0=1$. $(1,1)$: $u\in X\cap Z$, $u\in X\cup Z$; producto $1$, unión $1+1-1=1$. Las fórmulas coinciden con las funciones características en los cuatro casos.
  \item Complemento: $u\in\overline{X}\iff\chi_X(u)=0\iff1-\chi_X(u)=1$.
  \item Si $U$ es finito, $\abs{X}=\sum_{u\in U}\chi_X(u)$. Sumando la fórmula de la unión sobre $U$: $\abs{X\cup Z}=\abs{X}+\abs{Z}-\abs{X\cap Z}$.
\end{enumerate}
\resultado las tres identidades y el principio de inclusión-exclusión para dos conjuntos.
\comprobacion con los datos del \cref{eje:1-demorgan-num}: $\abs{A\cup B}=5+5-2=8$, y $A\cup B=\{1,2,3,4,5,6,8,10\}$ tiene 8 elementos.
\end{solucion}

\begin{solucion}{eje:2-biendef}
\planteamiento si $\tfrac ab=\tfrac cd$, comprobar si las imágenes coinciden; si no, dar dos representaciones con imágenes distintas.
\pasos
\begin{enumerate}[label=(\alph*)]
  \item $\tfrac{a^{2}}{b^{2}}=\bigl(\tfrac ab\bigr)^{2}$ depende solo del número $\tfrac ab$: \textbf{bien definida}.
  \item $\tfrac{a+b}{b}=\tfrac ab+1$ depende solo de $\tfrac ab$: \textbf{bien definida}.
  \item $\tfrac11=\tfrac22$, pero $h$ daría $\tfrac12$ y $\tfrac23$: \textbf{no está bien definida}.
\end{enumerate}
\resultado (a) y (b) son funciones; (c) no lo es.
\comprobacion en (b), con $\tfrac12=\tfrac24$: $\tfrac{3}{2}=\tfrac{6}{4}$.
\end{solucion}

\begin{solucion}{eje:2-argmax}
\datos $g(z)=$ menor índice $i$ con $z_i=\max(z_1,z_2,z_3)$.
\planteamiento evaluar y traducir la regla de desempate a desigualdades.
\pasos
\begin{enumerate}[label=\arabic*.]
  \item $g(1,3,2)=2$; $g(2,2,1)=1$ (empate entre 1 y 2, gana el menor); $g(0,0,0)=1$.
  \item $g^{-1}(\{1\})=\{z: z_1\ge z_2,\ z_1\ge z_3\}$: el índice~1 gana cuando es máximo, aunque haya empate.
  \item $g^{-1}(\{3\})=\{z: z_3>z_1,\ z_3>z_2\}$: el~3 solo gana si es estrictamente mayor.
  \item $g$ es una función de dominio $\R^{3}$: todo $z$ está en exactamente una fibra. La unión es $\R^{3}$ y las fibras son disjuntas.
\end{enumerate}
\resultado $2$, $1$, $1$; las tres fibras forman una partición de $\R^{3}$.
\comprobacion $(2,2,1)$ cumple $z_1\ge z_2$ y $z_1\ge z_3$, luego está en $g^{-1}(\{1\})$, como se obtuvo en el paso~1.
\end{solucion}
% ---------------------------------------------------------------------
\section{Tema 3}

\begin{solucion}{eje:3-tablas}
\planteamiento con dominio y codominio de 3 elementos, basta mirar si la imagen tiene 3 elementos (\cref{prop:t3-finitos}(ii)).
\pasos
\begin{enumerate}[label=(\alph*)]
  \item Imagen $\{\cad{a},\cad{b},\cad{c}\}$: \textbf{biyectiva}.
  \item $\cad{a}$ tiene dos preimágenes y $\cad{b}$ ninguna: \textbf{ni inyectiva ni sobreyectiva}.
  \item Imagen $\{\cad{c},\cad{a},\cad{b}\}$: \textbf{biyectiva}.
  \item Constante: \textbf{ni inyectiva ni sobreyectiva}.
  \item Ninguna puede ser inyectiva sin ser sobreyectiva: entre conjuntos finitos del mismo cardinal ambas propiedades son equivalentes.
\end{enumerate}
\resultado (a) y (c) biyectivas; (b) y (d), ninguna de las dos propiedades.
\comprobacion por la \cref{prop:t3-contar-iny} hay $3!=6$ biyecciones de $\{1,2,3\}$ en $\{\cad{a},\cad{b},\cad{c}\}$; (a) y (c) son dos de ellas.
\end{solucion}

\begin{solucion}{eje:3-error}
\planteamiento comparar la implicación probada con la definición de inyectividad.
\pasos
\begin{enumerate}[label=\arabic*.]
  \item El argumento prueba «$a=b\Rightarrow f(a)=f(b)$», que es cierto para cualquier función (dice que $f$ está bien definida).
  \item La inyectividad es la implicación recíproca: «$f(a)=f(b)\Rightarrow a=b$».
  \item Esta es falsa para $n^{2}$ en $\Z$: $f(-1)=1=f(1)$ con $-1\neq1$.
\end{enumerate}
\resultado el error es probar la implicación en el sentido equivocado; $f$ \textbf{no} es inyectiva.
\comprobacion en $\N$ la misma fórmula sí es inyectiva (\cref{ej:t3-cuadrado}, con $[0,\infty)$).
\end{solucion}

\begin{solucion}{eje:3-clasificar}
\planteamiento para cada función, demostración breve o contraejemplo.
\pasos
\begin{enumerate}[label=(\alph*)]
  \item $n+3=m+3\Rightarrow n=m$; $f(k-3)=k$. \textbf{Biyectiva}.
  \item Inyectiva ($2n=2m\Rightarrow n=m$); $1$ no es par: \textbf{inyectiva, no sobreyectiva}.
  \item Estrictamente creciente (\cref{eje:3-creciente}), luego inyectiva; $2$ no es cubo de un entero ($1^{3}<2<2^{3}$): \textbf{inyectiva, no sobreyectiva}.
  \item $n+1=m+1\Rightarrow n=m$; $0$ no es $n+1$ con $n\ge0$: \textbf{inyectiva, no sobreyectiva}.
  \item $f(0)=f(1)=0$: no inyectiva. Es continua con límites $-\infty$ y $+\infty$, luego toma todo valor real (teorema de Bolzano): \textbf{sobreyectiva, no inyectiva}.
  \item Estrictamente creciente, luego inyectiva; $e^{x}>0$, así que $-1$ no se alcanza: \textbf{inyectiva, no sobreyectiva}.
  \item $\lfloor0\rfloor=\lfloor0{,}5\rfloor$: no inyectiva; $\lfloor k\rfloor=k$ para $k\in\Z$: \textbf{sobreyectiva, no inyectiva}.
  \item $\abs{-1}=\abs{1}$: no inyectiva; $\abs{k}=k$ para $k\in\N$: \textbf{sobreyectiva, no inyectiva}.
\end{enumerate}
\resultado biyectiva: (a); solo inyectivas: (b), (c), (d), (f); solo sobreyectivas: (e), (g), (h).
\comprobacion (e): $x^{3}-x=x(x-1)(x+1)$ se anula en $-1$, $0$ y $1$, tres preimágenes del~0.
\end{solucion}

\begin{solucion}{eje:3-5n2}
\datos $f(n)=5n+2$.
\planteamiento inyectividad directa; para la no sobreyectividad, un valor sin preimagen entera.
\pasos
\begin{enumerate}[label=\arabic*.]
  \item Si $5n+2=5m+2$, entonces $5n=5m$ y $n=m$: inyectiva.
  \item $f(n)=0\iff n=-\tfrac25\notin\Z$: el $0$ no tiene preimagen. (En general, $f(n)$ deja resto~2 al dividir entre~5.)
  \item En $\Q$: la inyectividad se prueba igual; dado $y\in\Q$, $x=\tfrac{y-2}{5}\in\Q$ y $f(x)=y$. Es biyectiva.
\end{enumerate}
\resultado en $\Z$, inyectiva y no sobreyectiva; en $\Q$, biyectiva.
\comprobacion $f(-1)=-3$ y $f(0)=2$: los valores saltan de 5 en 5, y $0$ queda entre ellos.
\end{solucion}

\begin{solucion}{eje:3-homografica}
\datos $f(x)=\tfrac{3x+1}{x-2}$, dominio $\R\setminus\{2\}$, codominio $\R\setminus\{3\}$.
\planteamiento resolver $f(x)=y$ y comprobar que la solución está en el dominio.
\pasos
\begin{enumerate}[label=\arabic*.]
  \item Si $\tfrac{3x+1}{x-2}=3$, entonces $3x+1=3x-6$, es decir, $1=-6$: imposible. Luego $f(x)\neq3$.
  \item Inyectiva: si $f(x)=f(x')$, multiplicando en cruz, $(3x+1)(x'-2)=(3x'+1)(x-2)$, es decir, $3xx'-6x+x'-2=3xx'-6x'+x-2$; simplificando, $7x'=7x$ y $x=x'$.
  \item Sobreyectiva: sea $y\neq3$. De $y(x-2)=3x+1$ se obtiene $x(y-3)=2y+1$ y $x=\tfrac{2y+1}{y-3}$, que existe porque $y\neq3$.
  \item $x\neq2$: si $\tfrac{2y+1}{y-3}=2$, entonces $2y+1=2y-6$, imposible.
  \item Verificación: $3x+1=\tfrac{6y+3+y-3}{y-3}=\tfrac{7y}{y-3}$ y $x-2=\tfrac{2y+1-2y+6}{y-3}=\tfrac{7}{y-3}$, luego $f(x)=y$.
\end{enumerate}
\resultado $f$ es biyectiva y la preimagen de $y$ es $\dfrac{2y+1}{y-3}$.
\comprobacion $f(5)=\tfrac{16}{3}$ y $\tfrac{2\cdot16/3+1}{16/3-3}=\tfrac{35/3}{7/3}=5$.
\end{solucion}

\begin{solucion}{eje:3-lineal-z2}
\datos $f(m,n)=(m+n,\ m+2n)$ en $\Z\times\Z$.
\planteamiento resolver el sistema $m+n=a$, $m+2n=b$.
\pasos
\begin{enumerate}[label=\arabic*.]
  \item Inyectiva: si $(m+n,m+2n)=(m'+n',m'+2n')$, restando las componentes, $n=n'$, y entonces $m=m'$.
  \item Sobreyectiva: dado $(a,b)\in\Z\times\Z$, restando las ecuaciones $n=b-a$ y $m=a-n=2a-b$, ambos enteros.
  \item Verificación: $f(2a-b,\ b-a)=\bigl(2a-b+b-a,\ 2a-b+2b-2a\bigr)=(a,b)$.
\end{enumerate}
\resultado $f$ es biyectiva. \hfill$\blacksquare$
\comprobacion $f(1,1)=(2,3)$ y $(2\cdot2-3,\ 3-2)=(1,1)$.
\end{solucion}

\begin{solucion}{eje:3-2m3n}
\datos $f(m,n)=2m+3n$; $g(m,n)=2m+4n$.
\planteamiento para la sobreyectividad basta obtener el~1 y escalar.
\pasos
\begin{enumerate}[label=\arabic*.]
  \item $2\cdot(-1)+3\cdot1=1$. Dado $k\in\Z$, $f(-k,k)=-2k+3k=k$: sobreyectiva.
  \item $f(3,0)=6=f(0,2)$ con $(3,0)\neq(0,2)$: no inyectiva.
  \item $g(m,n)=2(m+2n)$ es siempre par; $1$ no tiene preimagen: $g$ no es sobreyectiva.
\end{enumerate}
\resultado $f$ es sobreyectiva y no inyectiva; $g$ no es sobreyectiva.
\comprobacion $f(-5,5)=-10+15=5$.
\end{solucion}

\begin{solucion}{eje:3-NZ}
\datos $f(n)=n/2$ ($n$ par), $f(n)=-(n+1)/2$ ($n$ impar).
\planteamiento separar por paridad: los pares van a $\N$ y los impares a los negativos.
\pasos
\begin{enumerate}[label=\arabic*.]
  \item $f(0)=0$, $f(1)=-1$, $f(2)=1$, $f(3)=-2$, $f(4)=2$, $f(5)=-3$, $f(6)=3$.
  \item Inyectiva. Sea $f(n)=f(m)$. Si $n$ es par, $f(n)\ge0$; si es impar, $f(n)\le-1$. Por tanto $n$ y $m$ tienen la misma paridad. Si son pares, $n/2=m/2$ da $n=m$; si son impares, $-(n+1)/2=-(m+1)/2$ da $n=m$.
  \item Sobreyectiva. Sea $k\in\Z$. Si $k\ge0$, $n=2k\in\N$ es par y $f(n)=k$. Si $k<0$, $n=-2k-1\ge1$ es impar y $f(n)=-(-2k-1+1)/2=k$.
\end{enumerate}
\resultado $f$ es biyectiva; lista $\Z$ como $0,-1,1,-2,2,-3,3,\dots$ \hfill$\blacksquare$
\comprobacion $k=-3$: $n=5$ y $f(5)=-3$, como en el paso~1.
\end{solucion}

\begin{solucion}{eje:3-creciente}
\datos $f,g\colon\R\to\R$.
\planteamiento (a) factorizar $y^{3}-x^{3}$; (b) aplicar la hipótesis dos veces; (c) buscar una biyección que intercambie dos puntos.
\pasos
\begin{enumerate}[label=\arabic*.]
  \item $x^{2}$ no es creciente: $-1<0$ pero $(-1)^{2}>0^{2}$.
  \item $x^{3}$ sí: si $x<y$, $y^{3}-x^{3}=(y-x)(y^{2}+xy+x^{2})$ con $y-x>0$ y $y^{2}+xy+x^{2}=(y+\tfrac x2)^{2}+\tfrac34x^{2}>0$, pues solo se anularía con $x=y=0$.
  \item (b) Si $x<y$, entonces $f(x)<f(y)$ y, aplicando $g$, $g(f(x))<g(f(y))$.
  \item (c) Falso. Sea $h(0)=1$, $h(1)=0$ y $h(x)=x$ en los demás puntos. $h\circ h=\id$, luego $h$ es biyectiva (en particular inyectiva). No es creciente ($0<1$ y $h(0)>h(1)$) ni decreciente ($1<2$ y $h(1)=0<2=h(2)$).
\end{enumerate}
\resultado $x^{3}$ es estrictamente creciente y $x^{2}$ no; la composición de estrictamente crecientes lo es; existen funciones inyectivas no monótonas. \hfill$\blacksquare$
\comprobacion para funciones \emph{continuas} en un intervalo sí es cierto que inyectiva implica estrictamente monótona; el contraejemplo es discontinuo en $0$ y en~$1$.
\end{solucion}

\begin{solucion}{eje:3-cadenas}
\datos $\Sigma=\{\cad{a},\cad{b}\}$; $r$ invierte, $d(w)=ww$, $q$ suprime el último símbolo.
\planteamiento usar longitudes y comparar símbolo a símbolo.
\pasos
\begin{enumerate}[label=\arabic*.]
  \item (a) $r(r(w))=w$ para toda $w$. Si $r(w)=r(w')$, aplicando $r$, $w=w'$: inyectiva. Dada $w$, $r(r(w))=w$: sobreyectiva. \textbf{Biyectiva}.
  \item (b) Si $ww=w'w'$, las longitudes dan $2\abs{w}=2\abs{w'}$; los primeros $\abs{w}$ símbolos de ambos lados son $w$ y $w'$, luego $w=w'$: inyectiva. $\abs{d(w)}$ es par, así que \cad{a} no está en la imagen: \textbf{inyectiva, no sobreyectiva}.
  \item (c) $q(\cad{aa})=\cad{a}=q(\cad{ab})$: no inyectiva. Dada $w$, $q(w\cad{a})=w$: \textbf{sobreyectiva, no inyectiva}.
\end{enumerate}
\resultado $r$ biyectiva; $d$ inyectiva y no sobreyectiva; $q$ sobreyectiva y no inyectiva.
\comprobacion $r(\cad{aab})=\cad{baa}$ y $r(\cad{baa})=\cad{aab}$.
\end{solucion}

\begin{solucion}{eje:3-codigos}
\datos $c\colon\Sigma\to\{0,1\}^{k}$ inyectiva; extensión por concatenación.
\planteamiento (a) cortar la cadena de bits en bloques de $k$; (b) un contraejemplo explícito.
\pasos
\begin{enumerate}[label=\arabic*.]
  \item (a) Sean $x=x_1\cdots x_n$, $y=y_1\cdots y_m$ con $c^{*}(x)=c^{*}(y)$. Las longitudes son $kn=km$, luego $n=m$.
  \item El bloque $i$-ésimo de $k$ bits es $c(x_i)$ en un lado y $c(y_i)$ en el otro; como las cadenas coinciden, $c(x_i)=c(y_i)$, y por la inyectividad de $c$, $x_i=y_i$ para todo $i$. Luego $x=y$.
  \item (b) $0$, $01$ y $1$ son distintas: $c$ es inyectiva. Pero $c^{*}(\cad{xz})=0\cdot1=01=c^{*}(\cad{y})$ con $\cad{xz}\neq\cad{y}$.
  \item (c) Es el mismo fenómeno: la decodificación concatena tokens, y $(\cad{a},\cad{b})$ y $(\cad{ab})$ dan el texto «ab». Los códigos de longitud fija, como los bytes, evitan el problema; UTF-8 lo evita con códigos de longitud variable sin prefijos (módulo~12, §1.4).
\end{enumerate}
\resultado $c^{*}$ es inyectiva para códigos de longitud fija; con longitud variable puede no serlo. \hfill$\blacksquare$
\comprobacion el código de (b) no es libre de prefijos: $c(\cad{x})=0$ es prefijo de $c(\cad{y})=01$.
\end{solucion}

\begin{solucion}{eje:3-contar}
\planteamiento \cref{prop:t2-contar,prop:t3-contar-iny}; para las sobreyecciones de (b), recuento directo.
\pasos
\begin{enumerate}[label=\arabic*.]
  \item (a) Funciones $4^{3}=64$; inyectivas $4\cdot3\cdot2=24$; sobreyectivas $0$ (pues $3<4$); biyectivas $0$.
  \item (b) Funciones $3^{4}=81$; inyectivas $0$ (pues $4>3$); biyectivas $0$.
  \item Sobreyecciones de (b): las fibras tienen tamaños $2,1,1$. Se elige el valor con dos preimágenes (3 formas), sus dos preimágenes ($\binom42=6$ formas) y se reparten los dos elementos restantes entre los otros dos valores (2 formas): $3\cdot6\cdot2=36$.
  \item (c) Funciones $3^{3}=27$; inyectivas, sobreyectivas y biyectivas, $3!=6$ (\cref{prop:t3-finitos}(ii)).
\end{enumerate}
\resultado (a) $64,24,0,0$; (b) $81,0,36,0$; (c) $27,6,6,6$.
\comprobacion la enumeración exhaustiva de las $81$ funciones de (b) da exactamente $36$ sobreyectivas.
\end{solucion}

\begin{solucion}{eje:3-finitos}
\datos $\abs{A}=\abs{B}=n$; $f$ sobreyectiva.
\planteamiento sumar los tamaños de las fibras.
\pasos
\begin{enumerate}[label=\arabic*.]
  \item $n=\abs{A}=\sum_{b\in B}\abs{f^{-1}(\{b\})}$, suma de $n$ términos.
  \item Por la sobreyectividad, cada término es $\ge1$. Si alguno fuese $\ge2$, la suma sería $\ge n+1$. Luego todos valen~1.
  \item Fibras de un elemento significan inyectividad.
  \item Contraejemplo infinito: $f\colon\N\to\N$, $f(n)=\lfloor n/2\rfloor$, es sobreyectiva ($f(2k)=k$) y no inyectiva ($f(0)=f(1)$).
\end{enumerate}
\resultado entre conjuntos finitos del mismo cardinal, sobreyectiva implica inyectiva; para conjuntos infinitos, no. \hfill$\blacksquare$
\comprobacion en el contraejemplo, $\N$ y $\N$ tienen «el mismo cardinal», pero todas las fibras tienen dos elementos.
\end{solucion}

\begin{solucion}{eje:3-softsign}
\datos $f(x)=\tfrac{x}{1+\abs{x}}$.
\planteamiento $f(x)$ tiene el signo de $x$; se trabaja por separado en $x\ge0$ y $x<0$.
\pasos
\begin{enumerate}[label=\arabic*.]
  \item (a) $\abs{f(x)}=\tfrac{\abs{x}}{1+\abs{x}}<1$ porque $\abs{x}<1+\abs{x}$.
  \item (b) Si $f(x)=f(x')$, $x$ y $x'$ tienen el mismo signo (o ambos son~0). Si $x,x'\ge0$: $x(1+x')=x'(1+x)$, es decir, $x+xx'=x'+xx'$ y $x=x'$. Si $x,x'<0$: $x(1-x')=x'(1-x)$, es decir, $x-xx'=x'-xx'$ y $x=x'$.
  \item (c) Sea $y\in[0,1)$ y $x=\tfrac{y}{1-y}\ge0$: $f(x)=\tfrac{y/(1-y)}{(1-y+y)/(1-y)}=y$. Sea $y\in(-1,0)$ y $x=\tfrac{y}{1+y}<0$: $f(x)=\tfrac{x}{1-x}=\tfrac{y/(1+y)}{(1+y-y)/(1+y)}=y$.
\end{enumerate}
\resultado $f$ es una biyección $\R\to(-1,1)$ con preimagen $x=\dfrac{y}{1-\abs{y}}$. \hfill$\blacksquare$
\comprobacion $y=\tfrac12$: $x=1$ y $f(1)=\tfrac12$; $y=-\tfrac34$: $x=-3$ y $f(-3)=-\tfrac34$.
\end{solucion}

\begin{solucion}{eje:3-tanh}
\datos $\tanh x=\tfrac{e^{x}-e^{-x}}{e^{x}+e^{-x}}=\tfrac{e^{2x}-1}{e^{2x}+1}$.
\planteamiento despejar $e^{2x}$.
\pasos
\begin{enumerate}[label=\arabic*.]
  \item $y=\tfrac{t-1}{t+1}$ con $t=e^{2x}>0$; entonces $y(t+1)=t-1$ y $t=\tfrac{1+y}{1-y}$.
  \item Para $y\in(-1,1)$, $\tfrac{1+y}{1-y}>0$, luego $x=\tfrac12\ln\tfrac{1+y}{1-y}$ existe y es la \emph{única} solución (cada paso es reversible y $\ln$ es inyectiva). Esto prueba a la vez sobreyectividad e inyectividad.
  \item Para $\abs{y}\ge1$ no hay solución: $\tanh x\in(-1,1)$, de modo que el codominio es el correcto.
  \item $y=0{,}6$: $x=\tfrac12\ln\tfrac{1{,}6}{0{,}4}=\tfrac12\ln4=\ln2\approx0{,}6931$.
\end{enumerate}
\resultado $\tanh$ es biyectiva con inversa $\tfrac12\ln\tfrac{1+y}{1-y}$; la preimagen de $0{,}6$ es $\ln2\approx0{,}6931$.
\comprobacion $\tanh(\ln2)=\tfrac{2-1/2}{2+1/2}=\tfrac{1{,}5}{2{,}5}=0{,}6$.
\end{solucion}

\begin{solucion}{eje:3-fibras}
\planteamiento traducir las definiciones a fibras.
\pasos
\begin{enumerate}[label=\arabic*.]
  \item Inyectiva: cada fibra tiene a lo sumo un elemento. Sobreyectiva: cada fibra tiene al menos uno. Biyectiva: cada fibra tiene exactamente uno.
  \item Cada $a\in A$ está en una única fibra, la de $f(a)$: las fibras forman una partición de $A$ (algunas pueden ser vacías). El cardinal de una unión disjunta es la suma de los cardinales.
\end{enumerate}
\resultado $\abs{A}=\sum_b\abs{f^{-1}(\{b\})}$, y las tres propiedades son condiciones sobre el tamaño de las fibras.
\comprobacion en el \cref{eje:2-tabla}, las fibras de $1,2,3,4,5$ tienen tamaños $1,3,0,1,0$, que suman $5=\abs{A}$.
\end{solucion}

% ---------------------------------------------------------------------
\section{Tema 4}

\begin{solucion}{eje:4-calc-comp}
\datos $f(x)=x^{2}+1$, $g(x)=2x-3$.
\planteamiento sustituir la función interior en la exterior.
\pasos
\begin{enumerate}[label=\arabic*.]
  \item $(g\circ f)(x)=2(x^{2}+1)-3=2x^{2}-1$.
  \item $(f\circ g)(x)=(2x-3)^{2}+1=4x^{2}-12x+10$.
  \item $(f\circ f)(x)=(x^{2}+1)^{2}+1=x^{4}+2x^{2}+2$.
  \item $(g\circ g)(x)=2(2x-3)-3=4x-9$.
  \item En $x=1$: $1$, $2$, $5$ y $-5$.
\end{enumerate}
\resultado $2x^{2}-1$, $4x^{2}-12x+10$, $x^{4}+2x^{2}+2$, $4x-9$; en $x=1$, $1,2,5,-5$.
\comprobacion directamente: $f(1)=2$, $g(2)=1$; $g(1)=-1$, $f(-1)=2$; $f(2)=5$; $g(-1)=-5$.
\end{solucion}

\begin{solucion}{eje:4-vf}
\planteamiento contraejemplos o resultados del tema.
\pasos
\begin{enumerate}[label=(\alph*)]
  \item \textbf{Falsa}: \cref{ej:t4-comp}.
  \item \textbf{Falsa} en general: $f^{-1}\circ f=\id_A$ y $f\circ f^{-1}=\id_B$, que solo son iguales si $A=B$.
  \item \textbf{Falsa}: con $f(x)=2x$, $f^{-1}(2)=1$ pero $1/f(2)=\tfrac14$.
  \item \textbf{Verdadera}: \cref{prop:t4-asoc}.
\end{enumerate}
\resultado F, F, F, V.
\comprobacion en (c) la igualdad sí se da para $f(x)=1/x$, que es su propia inversa; un caso particular no prueba la regla general.
\end{solucion}

\begin{solucion}{eje:4-tablas}
\datos $f=(1\mapsto2,\ 2\mapsto3,\ 3\mapsto1)$, $g=(1\mapsto1,\ 2\mapsto1,\ 3\mapsto3)$.
\planteamiento evaluar elemento a elemento; $(g\circ f)(i)=g(f(i))$.
\pasos
\begin{enumerate}[label=\arabic*.]
  \item $g\circ f$: $1\mapsto g(2)=1$, $2\mapsto g(3)=3$, $3\mapsto g(1)=1$.
  \item $f\circ g$: $1\mapsto f(1)=2$, $2\mapsto f(1)=2$, $3\mapsto f(3)=1$.
  \item $f\circ f$: $1\mapsto3$, $2\mapsto1$, $3\mapsto2$; $f\circ f\circ f=\id$.
  \item $f$ es biyectiva: tiene inversa. $g$ y $g\circ f$ no son inyectivas: no la tienen.
  \item $f^{-1}$: $1\mapsto3$, $2\mapsto1$, $3\mapsto2$, que coincide con $f\circ f$ porque $f\circ(f\circ f)=\id$.
\end{enumerate}
\resultado $g\circ f=(1,3,1)$, $f\circ g=(2,2,1)$, $f\circ f=(3,1,2)$, $f^{3}=\id$; solo $f$ es invertible, con $f^{-1}=f\circ f$.
\comprobacion $f(f^{-1}(1))=f(3)=1$, $f(f^{-1}(2))=f(1)=2$, $f(f^{-1}(3))=f(2)=3$.
\end{solucion}

\begin{solucion}{eje:4-inversas-reales}
\planteamiento despejar $x$ y verificar ambas composiciones.
\pasos
\begin{enumerate}[label=(\alph*)]
  \item $y=2x-7\iff x=\tfrac{y+7}{2}$: $f^{-1}(y)=\tfrac{y+7}{2}$. $f(f^{-1}(y))=y+7-7=y$; $f^{-1}(f(x))=\tfrac{2x}{2}=x$.
  \item $y=\ln(3x)\iff x=\tfrac{e^{y}}{3}>0$: $g^{-1}(y)=\tfrac{e^{y}}{3}$. $g(g^{-1}(y))=\ln e^{y}=y$; $g^{-1}(g(x))=\tfrac{3x}{3}=x$.
  \item $y=\sqrt{x-1}$ con $y\ge0$ $\iff x=y^{2}+1\ge1$: $h^{-1}(y)=y^{2}+1$. $h(h^{-1}(y))=\sqrt{y^{2}}=y$ (pues $y\ge0$); $h^{-1}(h(x))=x-1+1=x$.
\end{enumerate}
\resultado $f^{-1}(y)=\tfrac{y+7}{2}$, $g^{-1}(y)=\tfrac{e^{y}}{3}$, $h^{-1}(y)=y^{2}+1$. En cada caso la existencia y unicidad de la solución prueba la biyectividad.
\comprobacion $f(3)=-1$ y $f^{-1}(-1)=3$; $h(5)=2$ y $h^{-1}(2)=5$.
\end{solucion}

\begin{solucion}{eje:4-inversa-z2}
\datos $f(m,n)=(m+n,\ m+2n)$.
\planteamiento la solución del sistema obtenida en el \cref{eje:3-lineal-z2}.
\pasos
\begin{enumerate}[label=\arabic*.]
  \item $f^{-1}(a,b)=(2a-b,\ b-a)$.
  \item $f^{-1}(f(m,n))=\bigl(2(m+n)-(m+2n),\ (m+2n)-(m+n)\bigr)=(m,n)$.
  \item $f(f^{-1}(a,b))=(a,b)$, calculado en el \cref{eje:3-lineal-z2}.
\end{enumerate}
\resultado $f^{-1}(a,b)=(2a-b,\ b-a)$.
\comprobacion $f(3,-1)=(2,1)$ y $f^{-1}(2,1)=(3,-1)$.
\end{solucion}

\begin{solucion}{eje:4-comp-sobre}
\datos $f\colon A\to B$, $g\colon B\to C$ sobreyectivas.
\planteamiento dado $c\in C$, buscar su preimagen en dos etapas.
\pasos
\begin{enumerate}[label=\arabic*.]
  \item Sea $c\in C$. Como $g$ es sobreyectiva, existe $b\in B$ con $g(b)=c$.
  \item Como $f$ es sobreyectiva, existe $a\in A$ con $f(a)=b$.
  \item Entonces $(g\circ f)(a)=g(b)=c$: $g\circ f$ es sobreyectiva.
  \item Si $f$ y $g$ son biyectivas, $g\circ f$ es inyectiva por el \cref{teo:t4-comp}(i) y sobreyectiva por lo anterior.
\end{enumerate}
\resultado la composición de sobreyecciones es sobreyectiva y la de biyecciones es biyectiva. \hfill$\blacksquare$
\comprobacion $f(x)=x+1$ y $g(x)=2x$ son biyecciones de $\R$; $g\circ f=2x+2$ también, con inversa $\tfrac y2-1=f^{-1}(g^{-1}(y))$.
\end{solucion}

\begin{solucion}{eje:4-cambridge}
\datos $g\colon A\to B$, $f\colon B\to C$; se estudia $f\circ g$.
\planteamiento la \cref{prop:t4-reciproco} da las afirmaciones ciertas; un único ejemplo refuta las otras dos.
\pasos
\begin{enumerate}[label=\arabic*.]
  \item Si $f\circ g$ es inyectiva, $g$ (la función interior) es inyectiva: \cref{prop:t4-reciproco}(i).
  \item Si $f\circ g$ es sobreyectiva, $f$ (la función exterior) es sobreyectiva: \cref{prop:t4-reciproco}(ii).
  \item Ejemplo: $A=\{1\}$, $B=\{1,2\}$, $C=\{1\}$, $g(1)=1$, $f(1)=f(2)=1$. Entonces $f\circ g\colon\{1\}\to\{1\}$ es biyectiva.
  \item En él, $f$ no es inyectiva ($f(1)=f(2)$) y $g$ no es sobreyectiva ($2\notin g(A)$).
\end{enumerate}
\resultado (a) $g$ debe ser inyectiva, $f$ no; (b) $f$ debe ser sobreyectiva, $g$ no. \hfill$\blacksquare$
\comprobacion el ejemplo es el de un «codificador» $g$ seguido de un «decodificador» $f$ con $f\circ g=\id$: el codificador es inyectivo y el decodificador sobreyectivo, como en un tokenizador.
\end{solucion}

\begin{solucion}{eje:4-iteradas}
\datos $f(n)=2n-1$; $f^{3}=f\circ f\circ f$.
\planteamiento (a) inducción; (b) y (c) agrupar $f^{3}$ de dos maneras y aplicar la \cref{prop:t4-reciproco}.
\pasos
\begin{enumerate}[label=\arabic*.]
  \item $f^{2}(n)=2(2n-1)-1=4n-3$; $f^{3}(n)=2(4n-3)-1=8n-7$.
  \item Conjetura: $f^{k}(n)=2^{k}n-2^{k}+1$. Para $k=1$: $2n-1$. Si vale para $k$: $f^{k+1}(n)=f(f^{k}(n))=2(2^{k}n-2^{k}+1)-1=2^{k+1}n-2^{k+1}+1$.
  \item (b) $f^{3}=f\circ(f\circ f)$; si es sobreyectiva, la función exterior $f$ es sobreyectiva.
  \item (c) $f^{3}=(f\circ f)\circ f$; si es inyectiva, la función interior $f$ es inyectiva.
\end{enumerate}
\resultado $f^{k}(n)=2^{k}(n-1)+1$; y $f^{3}$ sobreyectiva (inyectiva) implica $f$ sobreyectiva (inyectiva). \hfill$\blacksquare$
\comprobacion $f^{3}(2)=f(f(3))=f(5)=9$ y $8\cdot2-7=9$.
\end{solucion}

\begin{solucion}{eje:4-inv-comp}
\datos $f\colon A\to B$, $g\colon B\to C$ biyectivas.
\planteamiento comprobar que $h=f^{-1}\circ g^{-1}$ cumple las dos condiciones de inversa de $g\circ f$ y usar la unicidad (\cref{teo:t4-inversa}).
\pasos
\begin{enumerate}[label=\arabic*.]
  \item $h\circ(g\circ f)=f^{-1}\circ(g^{-1}\circ g)\circ f=f^{-1}\circ\id_B\circ f=f^{-1}\circ f=\id_A$, por la asociatividad.
  \item $(g\circ f)\circ h=g\circ(f\circ f^{-1})\circ g^{-1}=g\circ\id_B\circ g^{-1}=\id_C$.
  \item Por la unicidad de la inversa, $(g\circ f)^{-1}=f^{-1}\circ g^{-1}$.
\end{enumerate}
\resultado $(g\circ f)^{-1}=f^{-1}\circ g^{-1}$. \hfill$\blacksquare$
\comprobacion con $f(x)=x+1$, $g(x)=2x$: $(g\circ f)(x)=2x+2$, cuya inversa es $\tfrac y2-1$, y $f^{-1}(g^{-1}(y))=\tfrac y2-1$; en cambio $g^{-1}(f^{-1}(y))=\tfrac{y-1}{2}$ es distinta.
\end{solucion}

\begin{solucion}{eje:4-derecha}
\datos $f\colon A\to B$.
\planteamiento ($\Rightarrow$) \cref{prop:t4-reciproco}(ii); ($\Leftarrow$) elegir una preimagen de cada elemento.
\pasos
\begin{enumerate}[label=\arabic*.]
  \item Si $f\circ s=\id_B$, que es sobreyectiva, la función exterior $f$ es sobreyectiva.
  \item Si $f$ es sobreyectiva, cada fibra $f^{-1}(\{b\})$ es no vacía; se elige un elemento $s(b)$ de cada una. Entonces $f(s(b))=b$ para todo $b$, es decir, $f\circ s=\id_B$. (Si $B$ es infinito, esta elección simultánea es el axioma de elección.)
  \item No es única: para $f(n)=\lfloor n/2\rfloor$, tanto $s_1(m)=2m$ como $s_2(m)=2m+1$ cumplen $f(s_i(m))=m$.
\end{enumerate}
\resultado $f$ tiene inversa por la derecha si y solo si es sobreyectiva; en general hay muchas, una por cada forma de elegir preimágenes. \hfill$\blacksquare$
\comprobacion $f(s_2(3))=\lfloor7/2\rfloor=3$.
\end{solucion}

\begin{solucion}{eje:4-involucion}
\datos $f\colon A\to A$ con $f\circ f=\id_A$.
\planteamiento la condición dice que $f$ es inversa de sí misma.
\pasos
\begin{enumerate}[label=\arabic*.]
  \item (a) Con $g=f$, las dos condiciones $g\circ f=\id_A$ y $f\circ g=\id_A$ son la misma hipótesis. Por el \cref{teo:t4-inversa}, $f$ es biyectiva y $f^{-1}=f$.
  \item (b) $-(-x)=x$; $1/(1/x)=x$; $A\setminus(A\setminus X)=X$ (\cref{ej:t1-doble} con $B=X\subseteq A$); invertir dos veces una cadena devuelve la original.
  \item (c) $f(1)=2$, $f(2)=3$, $f(3)=1$: $f(f(1))=3\neq1$.
\end{enumerate}
\resultado toda involución es una biyección igual a su inversa; un ciclo de longitud~3 es biyectivo y no es una involución. \hfill$\blacksquare$
\comprobacion en (c), $f\circ f\circ f=\id$ (\cref{eje:4-tablas}), luego $f^{-1}=f\circ f\neq f$.
\end{solucion}

\begin{solucion}{eje:4-cancelar}
\datos $f\colon A\to B$.
\planteamiento ($\Rightarrow$) directo; ($\Leftarrow$) por contrarrecíproco, construyendo $g\neq h$ que coinciden en $f(A)$.
\pasos
\begin{enumerate}[label=\arabic*.]
  \item ($\Rightarrow$) Sea $f$ sobreyectiva y $g\circ f=h\circ f$. Dado $b\in B$, existe $a$ con $f(a)=b$, y $g(b)=g(f(a))=h(f(a))=h(b)$. Luego $g=h$.
  \item ($\Leftarrow$) Si $f$ no es sobreyectiva, existe $b_0\notin f(A)$. Sea $C=\{0,1\}$, $g$ la función constante~0 y $h=\chi_{\{b_0\}}$. Para todo $a$, $f(a)\neq b_0$, luego $g(f(a))=0=h(f(a))$; pero $g(b_0)=0\neq1=h(b_0)$.
  \item Análogo: $f$ es inyectiva si y solo si para todo $C$ y todas $g,h\colon C\to A$, $f\circ g=f\circ h$ implica $g=h$.
\end{enumerate}
\resultado las sobreyecciones se cancelan por la derecha y las inyecciones por la izquierda. \hfill$\blacksquare$
\comprobacion con $f\colon\{1\}\to\{1,2\}$, $f(1)=1$ (no sobreyectiva), $g=0$ y $h=\chi_{\{2\}}$ cumplen $g\circ f=h\circ f$ y $g\neq h$.
\end{solucion}

\begin{solucion}{eje:4-capas}
\datos $f_1(x,y)=(x+y,\ x-y)$, $f_2(u,v)=(2u,\ v+1)$.
\planteamiento sustituir y, para las inversas, resolver cada sistema; después \cref{prop:t4-inv-comp}.
\pasos
\begin{enumerate}[label=\arabic*.]
  \item $F(x,y)=f_2(x+y,\ x-y)=\bigl(2x+2y,\ x-y+1\bigr)$; $F(1,2)=(6,0)$.
  \item $u=x+y$, $v=x-y$ dan $x=\tfrac{u+v}{2}$, $y=\tfrac{u-v}{2}$: $f_1^{-1}(u,v)=\bigl(\tfrac{u+v}{2},\tfrac{u-v}{2}\bigr)$.
  \item $f_2^{-1}(p,q)=\bigl(\tfrac p2,\ q-1\bigr)$.
  \item $F^{-1}=f_1^{-1}\circ f_2^{-1}$: $F^{-1}(p,q)=\Bigl(\tfrac{p/2+q-1}{2},\ \tfrac{p/2-q+1}{2}\Bigr)$.
  \item $F^{-1}(6,0)=\bigl(\tfrac{3-1}{2},\tfrac{3+1}{2}\bigr)=(1,2)$.
\end{enumerate}
\resultado $F(x,y)=(2x+2y,\ x-y+1)$ y $F^{-1}(p,q)=\bigl(\tfrac p4+\tfrac{q-1}{2},\ \tfrac p4-\tfrac{q-1}{2}\bigr)$.
\comprobacion $F(F^{-1}(p,q))$: la primera componente es $2\bigl(\tfrac p4+\tfrac{q-1}2+\tfrac p4-\tfrac{q-1}2\bigr)=p$ y la segunda, $(q-1)+1=q$.
\end{solucion}

\begin{solucion}{eje:4-tokenizador}
\datos $\delta\circ\tau=\id_{\Sigma^{*}}$; $p=\tau\circ\delta$.
\planteamiento asociatividad y \cref{prop:t4-reciproco}.
\pasos
\begin{enumerate}[label=\arabic*.]
  \item (a) $p\circ p=\tau\circ(\delta\circ\tau)\circ\delta=\tau\circ\id\circ\delta=p$.
  \item (b, $\Rightarrow$) Si $p=\tau\circ\delta=\id_{\V^{*}}$, es inyectiva, luego la función interior $\delta$ es inyectiva.
  \item (b, $\Leftarrow$) Si $\delta$ es inyectiva, para todo $x\in\V^{*}$: $\delta(p(x))=(\delta\circ\tau)(\delta(x))=\delta(x)$, y por inyectividad $p(x)=x$.
  \item (c) $\delta((\cad{a},\cad{b}))=$ «ab»; la segmentación de máxima longitud toma \cad{ab} como un solo token: $p((\cad{a},\cad{b}))=(\cad{ab})$.
\end{enumerate}
\resultado $p$ es idempotente; $p=\id$ si y solo si $\delta$ es inyectiva; en el ejemplo, $p((\cad{a},\cad{b}))=(\cad{ab})\neq(\cad{a},\cad{b})$. \hfill$\blacksquare$
\comprobacion $p((\cad{ab}))=\tau(\text{«ab»})=(\cad{ab})$: aplicar $p$ otra vez no cambia nada, como predice~(a). La función $p$ lleva cada secuencia a la segmentación «canónica» del texto.
\end{solucion}
% ---------------------------------------------------------------------
\section{Tema 5}

\begin{solucion}{eje:5-propiedades}
\datos cuatro relaciones en $A=\{1,2,3\}$.
\planteamiento comprobar cada propiedad de la \cref{def:t5-prop}; para refutar basta un par concreto.
\pasos
\begin{enumerate}[label=\arabic*.]
  \item $R_1$: reflexiva (contiene $(1,1),(2,2),(3,3)$); simétrica; no antisimétrica ($1R_12$, $2R_11$, $1\neq2$); transitiva (todas las cadenas $aR_1bR_1c$ quedan dentro de $\{1,2\}$ o en $3R_13$). Es la equivalencia con clases $\{1,2\}$ y $\{3\}$.
  \item $R_2$: no reflexiva ($(1,1)\notin R_2$); no simétrica ($(2,1)\notin R_2$); antisimétrica (no hay pares $(a,b),(b,a)$ con $a\neq b$); transitiva (la única cadena es $1R_22R_23$, y $(1,3)\in R_2$).
  \item $R_3$ es la igualdad: reflexiva, simétrica, antisimétrica y transitiva.
  \item $R_4$: no reflexiva; simétrica; no antisimétrica; no transitiva ($1R_42$, $2R_41$, pero $(1,1)\notin R_4$).
\end{enumerate}
\resultado $R_1$: R, S, T; $R_2$: A, T; $R_3$: R, S, A, T; $R_4$: S.
\comprobacion $R_2$ es el orden estricto $<$ en $\{1,2,3\}$, que es antisimétrico y transitivo, como en la tabla del tema.
\end{solucion}

\begin{solucion}{eje:5-enteros}
\planteamiento demostración o contraejemplo para cada propiedad.
\pasos
\begin{enumerate}[label=(\alph*)]
  \item $a+b$ par $\iff$ $a$ y $b$ tienen la misma paridad. Reflexiva ($2a$ par), simétrica, transitiva (misma paridad es transitiva); no antisimétrica ($0R2$, $2R0$). \textbf{Equivalencia}, clases: pares e impares.
  \item $ab\ge0$. Reflexiva ($a^{2}\ge0$), simétrica; no antisimétrica ($1R2$, $2R1$); no transitiva: $1R0$ y $0R(-1)$, pero $1\cdot(-1)<0$.
  \item $\abs{a-b}\le1$. Reflexiva, simétrica; no antisimétrica ($0R1$, $1R0$); no transitiva: $0R1$, $1R2$, pero $\abs{0-2}=2$.
\end{enumerate}
\resultado solo (a) es de equivalencia, con clases $2\Z$ y $2\Z+1$.
\comprobacion (a) es el núcleo de $n\mapsto n\bmod2$, de acuerdo con la \cref{prop:t5-nucleo}.
\end{solucion}

\begin{solucion}{eje:5-contar}
\datos $\abs{A}=n=3$; $\abs{A\times A}=n^{2}=9$.
\planteamiento una relación es un subconjunto de $A\times A$: se decide par a par.
\pasos
\begin{enumerate}[label=\arabic*.]
  \item Todas: $2^{n^{2}}=2^{9}=512$.
  \item Reflexivas: los $n$ pares diagonales están obligados; los $n^{2}-n=6$ restantes son libres: $2^{6}=64$.
  \item Simétricas: se decide cada par diagonal ($n$) y cada pareja $\{(a,b),(b,a)\}$ con $a<b$ ($\binom n2$) a la vez: $2^{n+\binom n2}=2^{n(n+1)/2}=2^{6}=64$.
  \item Reflexivas y simétricas: solo las parejas fuera de la diagonal son libres: $2^{\binom n2}=2^{3}=8$.
\end{enumerate}
\resultado $512$, $64$, $64$, $8$; en general $2^{n^{2}}$, $2^{n^{2}-n}$, $2^{n(n+1)/2}$, $2^{n(n-1)/2}$.
\comprobacion la enumeración de las $512$ relaciones confirma los cuatro números.
\end{solucion}

\begin{solucion}{eje:5-clases}
\planteamiento describir $[a]$ para un $a$ genérico.
\pasos
\begin{enumerate}[label=(\alph*)]
  \item $[r]=\{4k+r:k\in\Z\}$ para $r=0,1,2,3$; el cociente $\Z_4=\{[0],[1],[2],[3]\}$ tiene 4 elementos.
  \item La clase de $p\neq0$ es la circunferencia de centro el origen y radio $\norm{p}$; la del origen es $\{(0,0)\}$. El cociente se identifica con $[0,\infty)$ mediante $[p]\mapsto\norm{p}$.
  \item La clase de $w$ es $\Sigma^{\abs{w}}$, con $2^{\abs{w}}$ elementos; el cociente se identifica con $\N$.
\end{enumerate}
\resultado (a) 4 clases; (b) circunferencias concéntricas y el origen; (c) una clase de $2^{n}$ cadenas por cada $n\in\N$.
\comprobacion en los tres casos la relación es el núcleo de una función (resto entre~4, distancia al origen, longitud) y las clases son sus fibras (\cref{prop:t5-nucleo}).
\end{solucion}

\begin{solucion}{eje:5-particion}
\datos bloques $\{1,2\}$, $\{3\}$, $\{4,5,6\}$.
\planteamiento cada bloque de tamaño $k$ aporta todos sus $k^{2}$ pares ordenados.
\pasos
\begin{enumerate}[label=\arabic*.]
  \item $\{1,2\}$: $(1,1),(1,2),(2,1),(2,2)$.
  \item $\{3\}$: $(3,3)$.
  \item $\{4,5,6\}$: los 9 pares $(a,b)$ con $a,b\in\{4,5,6\}$.
  \item Total: $4+1+9=14$; en general, $k_1^{2}+\dots+k_r^{2}$.
\end{enumerate}
\resultado 14 pares; $\sum_i k_i^{2}$ en general.
\comprobacion los casos extremos son correctos: la igualdad (bloques unitarios) tiene $n$ pares y la relación total (un bloque) tiene $n^{2}$.
\end{solucion}

\begin{solucion}{eje:5-bell}
\planteamiento por el \cref{teo:t5-particion}, contar relaciones de equivalencia es contar particiones.
\pasos
\begin{enumerate}[label=\arabic*.]
  \item $n=3$. Tipo $\{3\}$: 1. Tipo $\{2,1\}$: se elige el elemento aislado, 3. Tipo $\{1,1,1\}$: 1. Total~5.
  \item $n=4$. Tipo $\{4\}$: 1. Tipo $\{3,1\}$: 4. Tipo $\{2,2\}$: el compañero del~1 se elige de 3 formas, 3. Tipo $\{2,1,1\}$: $\binom42=6$. Tipo $\{1,1,1,1\}$: 1. Total $1+4+3+6+1=15$.
\end{enumerate}
\resultado 5 relaciones de equivalencia en un conjunto de 3 elementos y 15 en uno de 4 (números de Bell $B_3$ y $B_4$).
\comprobacion la recursión de los números de Stirling de segunda especie da $S(4,\cdot)=1,7,6,1$, que suma 15; y $7=4+3$ coincide con las particiones en dos bloques.
\end{solucion}

\begin{solucion}{eje:5-modn}
\datos $n\ge1$; $a\equiv b\iff n\mid(a-b)$.
\planteamiento propiedades de la divisibilidad y división euclídea.
\pasos
\begin{enumerate}[label=\arabic*.]
  \item Reflexiva: $a-a=0=0\cdot n$.
  \item Simétrica: si $a-b=kn$, entonces $b-a=(-k)n$.
  \item Transitiva: si $a-b=kn$ y $b-c=ln$, entonces $a-c=(k+l)n$.
  \item Por la división euclídea, todo $a$ se escribe $a=qn+r$ con $0\le r<n$; entonces $a-r=qn$ y $a\in[r]$. Hay a lo sumo $n$ clases.
  \item Si $0\le r<r'<n$, entonces $0<r'-r<n$ no es múltiplo de $n$, y $[r]\neq[r']$. Hay exactamente $n$ clases.
\end{enumerate}
\resultado la congruencia módulo $n$ es una equivalencia con las $n$ clases $[0],\dots,[n-1]$. \hfill$\blacksquare$
\comprobacion $n=3$: $-4=(-2)\cdot3+2$, luego $-4\in[2]$, como en el \cref{ej:t5-modn}.
\end{solucion}

\begin{solucion}{eje:5-nucleo}
\datos $f\colon A\to B$; $a\sim_f a'\iff f(a)=f(a')$.
\planteamiento las propiedades de $\sim_f$ son las de la igualdad en $B$.
\pasos
\begin{enumerate}[label=\arabic*.]
  \item Reflexiva: $f(a)=f(a)$. Simétrica: $f(a)=f(a')\Rightarrow f(a')=f(a)$. Transitiva: $f(a)=f(a')$ y $f(a')=f(a'')$ dan $f(a)=f(a'')$.
  \item $[a]=\{x: f(x)=f(a)\}=f^{-1}(\{f(a)\})$: cada clase es una fibra no vacía.
  \item Recíprocamente, si $f^{-1}(\{b\})\neq\emptyset$, contiene algún $a$, y entonces $b=f(a)$ y la fibra es $[a]$.
  \item $f$ es inyectiva $\iff$ todas las fibras no vacías tienen un elemento $\iff$ todas las clases son unitarias.
\end{enumerate}
\resultado $\sim_f$ es una equivalencia cuyas clases son las fibras no vacías; $f$ es inyectiva si y solo si $\sim_f$ es la igualdad. \hfill$\blacksquare$
\comprobacion para $f(x)=x^{2}$ en $\R$, $[3]=\{-3,3\}=f^{-1}(\{9\})$.
\end{solucion}

\begin{solucion}{eje:5-divide}
\datos $aRb\iff a\mid b$ o $b\mid a$, en $\N_{\ge1}$.
\planteamiento comprobar las tres propiedades.
\pasos
\begin{enumerate}[label=\arabic*.]
  \item Reflexiva: $a\mid a$.
  \item Simétrica: la condición «$a\mid b$ o $b\mid a$» no cambia al intercambiar $a$ y $b$.
  \item No transitiva: $2R6$ (pues $2\mid6$) y $6R3$ (pues $3\mid6$), pero $2\nmid3$ y $3\nmid2$, luego no $2R3$.
\end{enumerate}
\resultado $R$ es reflexiva y simétrica, pero \textbf{no} es de equivalencia porque no es transitiva.
\comprobacion si $R$ fuese de equivalencia, como $1$ está relacionado con todo $n$, habría una sola clase y $2R3$; se ha visto que no.
\end{solucion}

\begin{solucion}{eje:5-mas3}
\datos $aRb\iff b=a+3$.
\planteamiento probar que toda equivalencia $S\supseteq R$ contiene a la congruencia módulo~3, y contar las equivalencias que la contienen.
\pasos
\begin{enumerate}[label=\arabic*.]
  \item Sea $S$ de equivalencia con $R\subseteq S$. Por transitividad, $a\,S\,(a+3k)$ para todo $k\ge0$ (inducción sobre $k$); por simetría, también para $k<0$. Luego $a\equiv b\ (\mathrm{mod}\ 3)$ implica $a\,S\,b$.
  \item Cada clase de $S$ es unión de clases módulo~3: si $x$ está en una clase de $S$ e $y\equiv x$, entonces $y\,S\,x$. Así $S$ queda determinada por una partición del conjunto $\{[0],[1],[2]\}$.
  \item Recíprocamente, cada partición de $\{[0],[1],[2]\}$ define (uniendo bloques) una partición de $\Z$ y por tanto una equivalencia que contiene a la congruencia módulo~3, y así a $R$.
  \item Hay $B_3=5$ particiones de un conjunto de 3 elementos (\cref{eje:5-bell}).
\end{enumerate}
\resultado exactamente \textbf{5} relaciones de equivalencia contienen a $R$: la congruencia módulo~3, las tres que unen dos de sus clases y la relación total. \hfill$\blacksquare$
\comprobacion la menor es la congruencia módulo~3 (3 clases) y la mayor es $\Z\times\Z$ (1 clase); las tres intermedias tienen 2 clases: $1+3+1=5$.
\end{solucion}

\begin{solucion}{eje:5-fracciones}
\datos $(a,b)\sim(c,d)\iff ad=bc$, con $b,d\neq0$.
\planteamiento la transitividad requiere cancelar un factor no nulo.
\pasos
\begin{enumerate}[label=\arabic*.]
  \item Reflexiva: $ab=ba$. Simétrica: $ad=bc\Rightarrow cb=da$.
  \item Transitiva: sean $ad=bc$ y $cf=de$ con $b,d,f\neq0$. Multiplicando la primera por $f$ y usando la segunda: $adf=bcf=bde$. Luego $d(af-be)=0$ y, como $d\neq0$, $af=be$, es decir, $(a,b)\sim(e,f)$.
  \item $[(1,2)]=\{(c,d): d=2c,\ d\neq0\}=\{(k,2k): k\in\Z\setminus\{0\}\}$.
  \item La aplicación $[(a,b)]\mapsto a/b$ está bien definida ($ad=bc\iff a/b=c/d$) y es biyectiva sobre $\Q$: es inyectiva por la misma equivalencia y sobreyectiva porque todo racional es $a/b$.
\end{enumerate}
\resultado $\sim$ es una equivalencia; cada clase es «una fracción con todas sus representaciones», y $X/{\sim}\cong\Q$. \hfill$\blacksquare$
\comprobacion $(1,2)\sim(-3,-6)$ porque $1\cdot(-6)=-6=2\cdot(-3)$.
\end{solucion}

\begin{solucion}{eje:5-orden-partes}
\planteamiento comprobar reflexiva, antisimétrica y transitiva; refutar con pares concretos.
\pasos
\begin{enumerate}[label=\arabic*.]
  \item (a) $X\subseteq X$; si $X\subseteq Y$ e $Y\subseteq X$, $X=Y$ (doble inclusión); transitividad: \cref{prop:t1-subconj}(ii). Si $a\neq b$ están en $A$, ni $\{a\}\subseteq\{b\}$ ni $\{b\}\subseteq\{a\}$: no es total.
  \item (b) Reflexiva: $a=1\cdot a$. Transitiva: $b=ka$, $c=lb$ dan $c=(lk)a$.
  \item Antisimétrica en $\N_{\ge1}$: si $b=ka$ y $a=lb$ con $k,l\ge1$ enteros, entonces $a=lka$, $lk=1$ y $k=l=1$, luego $a=b$.
  \item En $\Z$: $2\mid-2$ y $-2\mid2$, pero $2\neq-2$.
\end{enumerate}
\resultado $\subseteq$ y $\mid$ (en $\N_{\ge1}$) son órdenes parciales no totales; $\mid$ no es antisimétrica en $\Z$. \hfill$\blacksquare$
\comprobacion en $\N_{\ge1}$, $2$ y $3$ son incomparables por divisibilidad, igual que $\{a\}$ y $\{b\}$ por inclusión.
\end{solucion}

\begin{solucion}{eje:5-hasse}
\datos divisores de 12: $1,2,3,4,6,12$; cadenas sobre $\cad{a}<\cad{b}$.
\planteamiento dibujar solo las relaciones de cobertura; para (b), comparar símbolo a símbolo.
\pasos
\begin{enumerate}[label=\arabic*.]
  \item Coberturas: $1\prec2$, $1\prec3$, $2\prec4$, $2\prec6$, $3\prec6$, $4\prec12$, $6\prec12$. El diagrama es:
  \begin{center}
  \begin{tikzpicture}[font=\small,scale=0.8]
    \node (1) at (0,0) {1}; \node (2) at (-1,1) {2}; \node (3) at (1,1) {3};
    \node (4) at (-2,2) {4}; \node (6) at (0,2) {6}; \node (12) at (-1,3) {12};
    \draw (1)--(2) (1)--(3) (2)--(4) (2)--(6) (3)--(6) (4)--(12) (6)--(12);
  \end{tikzpicture}
  \end{center}
  \item $1$ divide a todos: es mínimo y único minimal. $12$ es múltiplo de todos: máximo y único maximal.
  \item (b) $\vacia$ es prefijo de todas, va primera. Empiezan por \cad{a}: \cad{a} (prefijo de las demás), y \cad{aab} antes que \cad{ab} porque en la segunda posición $\cad{a}<\cad{b}$. Después \cad{b} y \cad{ba}.
\end{enumerate}
\resultado (a) mínimo 1, máximo 12; (b) $\vacia<\cad{a}<\cad{aab}<\cad{ab}<\cad{b}<\cad{ba}$.
\comprobacion la ordenación de cadenas de un lenguaje de programación (comparación lexicográfica) produce la misma lista.
\end{solucion}

\begin{solucion}{eje:5-lex}
\datos $(a,b)\preceq(c,d)\iff a<c$, o $a=c$ y $b\le d$.
\planteamiento comprobar las cuatro propiedades distinguiendo según la comparación de las primeras componentes.
\pasos
\begin{enumerate}[label=\arabic*.]
  \item Reflexiva: $a=a$ y $b\le b$.
  \item Antisimétrica: si $(a,b)\preceq(c,d)$ y $(c,d)\preceq(a,b)$, no puede ser $a<c$ (la segunda exigiría $c<a$ o $c=a$), ni $c<a$; luego $a=c$, y entonces $b\le d\le b$, así que $b=d$.
  \item Transitiva: sean $(a,b)\preceq(c,d)\preceq(e,f)$. Si $a<c$ o $c<e$, entonces $a\le c\le e$ con alguna desigualdad estricta, luego $a<e$. Si $a=c=e$, entonces $b\le d\le f$.
  \item Total: si $a<c$ o $c<a$, los pares son comparables; si $a=c$, se comparan $b$ y $d$, que siempre son comparables en $\N$.
  \item Orden producto: $(0,1)$ y $(1,0)$ no son comparables ($0\le1$ pero $1\not\le0$), luego no es total.
\end{enumerate}
\resultado el orden lexicográfico es total; el orden producto es parcial pero no total. \hfill$\blacksquare$
\comprobacion $(1,5)\preceq(2,0)$ en el orden lexicográfico, aunque $5>0$: manda la primera componente.
\end{solucion}

\begin{solucion}{eje:5-error}
\planteamiento identificar la hipótesis oculta del argumento.
\pasos
\begin{enumerate}[label=\arabic*.]
  \item El argumento empieza con «si $a\,R\,b$»: solo prueba $a\,R\,a$ para los $a$ que están relacionados con \emph{algún} $b$.
  \item Contraejemplo en $\{1,2\}$: $R=\{(1,1)\}$. Es simétrica y transitiva (trivialmente), pero no reflexiva, porque $(2,2)\notin R$; el~2 no está relacionado con nadie.
\end{enumerate}
\resultado el error es suponer que todo elemento está relacionado con alguno; $R=\{(1,1)\}$ es un contraejemplo.
\comprobacion la relación vacía en un conjunto no vacío es otro contraejemplo.
\end{solucion}

% ---------------------------------------------------------------------
\section{Tema 6}

\begin{solucion}{eje:6-vf}
\planteamiento remitir a los resultados del tema.
\pasos
\begin{enumerate}[label=(\alph*)]
  \item \textbf{Falsa} en general: $2\N\subsetneq\N$ y $\abs{2\N}=\abs{\N}$ (\cref{ej:t6-pares}). Es cierta si $B$ es finito.
  \item \textbf{Verdadera}: \cref{ej:t6-pares}.
  \item \textbf{Verdadera}: \cref{eje:6-Q}.
  \item \textbf{Falsa}: \cref{cor:t6-nonum}.
  \item \textbf{Verdadera}: ambos son infinitos numerables (\cref{prop:t6-numerables}); si $\alpha\colon\N\to\Z$ y $\pi\colon\N\times\N\to\N$ son biyecciones, $\alpha\circ\pi$ es una biyección $\N\times\N\to\Z$, y su inversa va de $\Z$ a $\N\times\N$.
\end{enumerate}
\resultado F, V, V, F, V.
\comprobacion (e) usa que la composición y la inversa de biyecciones son biyecciones (\cref{teo:t4-comp,prop:t4-inv-comp}).
\end{solucion}

\begin{solucion}{eje:6-textos}
\planteamiento comparar $\abs{\V^{*}}$ con $\abs{\N}$, $\abs{\R}$ y con un conjunto finito.
\pasos
\begin{enumerate}[label=\arabic*.]
  \item $\V^{*}$ es infinito numerable (\cref{prop:t6-numerables}): hay \emph{tantos} textos como naturales, ni más ni menos.
  \item $\R$ no es numerable: hay estrictamente menos textos que números reales.
  \item Los números en coma flotante de 32~bits son a lo sumo $2^{32}$, un conjunto finito; $\V^{*}$ es infinito: hay más textos.
\end{enumerate}
\resultado $\abs{\V^{*}}=\abs{\N}<\abs{\R}$, y $\V^{*}$ es mayor que cualquier conjunto finito.
\comprobacion en particular, no todo número real puede describirse con un texto finito: las descripciones forman un conjunto numerable.
\end{solucion}

\begin{solucion}{eje:6-meses}
\planteamiento \cref{teo:t6-palomar}; palomas = personas o calcetines; palomares = meses o colores.
\pasos
\begin{enumerate}[label=\arabic*.]
  \item (a) $13>12$: sí, dos personas comparten mes.
  \item (b) Con $24=2\cdot12$ personas puede haber exactamente dos por mes. Con $25>2\cdot12$, algún mes tiene al menos~3. Respuesta: 25.
  \item (c) Con 4 calcetines puede haber uno de cada color; con $5>4$, dos comparten color. Respuesta: 5.
\end{enumerate}
\resultado (a) sí; (b) 25 personas; (c) 5 calcetines.
\comprobacion $\lceil25/12\rceil=3$ y $\lceil5/4\rceil=2$.
\end{solucion}

\begin{solucion}{eje:6-bits}
\datos $K=5\cdot10^{4}$; $10^{6}$ palabras y $2^{16}=65\,536$ cubetas.
\planteamiento palomar en sus dos formas.
\pasos
\begin{enumerate}[label=\arabic*.]
  \item Con 15 bits hay $2^{15}=32\,768<50\,000$ identificadores: no existe una función inyectiva de los tokens en ellos.
  \item Con 16 bits hay $2^{16}=65\,536\ge50\,000$: basta. El mínimo es $\lceil\log_2 50\,000\rceil=16$.
  \item $15\cdot2^{16}=983\,040<10^{6}$, luego por la forma general (con $k=15$) alguna cubeta recibe al menos $16$ palabras.
\end{enumerate}
\resultado 16 bits como mínimo; alguna cubeta con al menos 16 palabras.
\comprobacion $10^{6}/65\,536\approx15{,}26$ y $\lceil15{,}26\rceil=16$.
\end{solucion}

\begin{solucion}{eje:6-restos}
\planteamiento elegir bien los palomares: restos en (a), parejas de consecutivos en (b).
\pasos
\begin{enumerate}[label=\arabic*.]
  \item (a) Palomares: los $n$ restos $0,\dots,n-1$ al dividir entre $n$. Con $n+1$ enteros, dos tienen el mismo resto, y su diferencia es múltiplo de $n$.
  \item (b) Palomares: las $n$ parejas $\{1,2\},\{3,4\},\dots,\{2n-1,2n\}$. Con $n+1$ números elegidos, dos caen en la misma pareja, y son consecutivos.
\end{enumerate}
\resultado ambas afirmaciones son ciertas. \hfill$\blacksquare$
\comprobacion el número $n+1$ es óptimo en (b): los $n$ impares $1,3,\dots,2n-1$ no contienen dos consecutivos.
\end{solucion}

\begin{solucion}{eje:6-puntos}
\datos cinco puntos $(x_i,y_i)\in\Z^{2}$.
\planteamiento palomares: las 4 clases de paridad $(x\bmod2,\ y\bmod2)$.
\pasos
\begin{enumerate}[label=\arabic*.]
  \item Hay 4 clases de paridad y 5 puntos: dos puntos, $(x_1,y_1)$ y $(x_2,y_2)$, tienen la misma clase.
  \item Entonces $x_1+x_2$ e $y_1+y_2$ son pares, y el punto medio $\bigl(\tfrac{x_1+x_2}{2},\tfrac{y_1+y_2}{2}\bigr)$ tiene coordenadas enteras.
  \item Con cuatro puntos es falso: $(0,0),(0,1),(1,0),(1,1)$ tienen paridades distintas y todos sus puntos medios tienen alguna coordenada $\tfrac12$.
\end{enumerate}
\resultado cierto para cinco puntos; falso en general para cuatro. \hfill$\blacksquare$
\comprobacion el punto medio de $(0,0)$ y $(1,1)$ es $(\tfrac12,\tfrac12)$.
\end{solucion}

\begin{solucion}{eje:6-compresion}
\datos $C\colon\{0,1\}^{*}\to\{0,1\}^{*}$ inyectiva.
\planteamiento (a) contar; (b) por reducción al absurdo con (a).
\pasos
\begin{enumerate}[label=\arabic*.]
  \item (a) $\abs{\{0,1\}^{n}}=2^{n}$ y $\abs{\{0,1\}^{<n}}=\sum_{k=0}^{n-1}2^{k}=2^{n}-1$ (suma geométrica). Como $2^{n}>2^{n}-1$, el palomar impide la inyectividad.
  \item (b) Supóngase que $\abs{C(w)}\le n-1$ para toda $w$ de longitud $n$. La restricción de $C$ a $\{0,1\}^{n}$ sería una función inyectiva en $\{0,1\}^{<n}$, en contra de~(a).
\end{enumerate}
\resultado ningún compresor sin pérdida acorta todas las cadenas de una longitud dada. \hfill$\blacksquare$
\comprobacion $n=2$: hay 4 cadenas de longitud~2 y solo $1+2=3$ de longitud menor.
\end{solucion}

\begin{solucion}{eje:6-pi}
\datos $\pi(m,n)=\tfrac{(m+n)(m+n+1)}{2}+n$.
\planteamiento evaluar; para invertir, buscar la diagonal $d$ con $\tfrac{d(d+1)}{2}\le17<\tfrac{(d+1)(d+2)}{2}$.
\pasos
\begin{enumerate}[label=\arabic*.]
  \item $\pi(2,3)=\tfrac{5\cdot6}{2}+3=18$; $\pi(0,4)=\tfrac{4\cdot5}{2}+4=14$.
  \item Números triangulares: $0,1,3,6,10,15,21$. Como $15\le17<21$, $d=5$, $n=17-15=2$ y $m=5-2=3$.
  \item Diagonal $d=3$: $\pi(3,0)=6$, $\pi(2,1)=7$, $\pi(1,2)=8$, $\pi(0,3)=9$, como en la \cref{fig:t6-cantor}.
\end{enumerate}
\resultado $\pi(2,3)=18$, $\pi(0,4)=14$, $\pi^{-1}(17)=(3,2)$.
\comprobacion $\pi(3,2)=\tfrac{5\cdot6}{2}+2=17$.
\end{solucion}

\begin{solucion}{eje:6-shortlex}
\datos $\Sigma=\{\cad{a},\cad{b}\}$.
\planteamiento por longitudes; dentro de cada longitud, orden lexicográfico (que coincide con el orden numérico binario si $\cad{a}=0$, $\cad{b}=1$).
\pasos
\begin{enumerate}[label=\arabic*.]
  \item $0{:}\ \vacia$; $1{:}\ \cad{a}$; $2{:}\ \cad{b}$; $3{:}\ \cad{aa}$; $4{:}\ \cad{ab}$; $5{:}\ \cad{ba}$; $6{:}\ \cad{bb}$; $7{:}\ \cad{aaa}$; $8{:}\ \cad{aab}$; $9{:}\ \cad{aba}$; $10{:}\ \cad{abb}$; $11{:}\ \cad{baa}$; $12{:}\ \cad{bab}$; $13{:}\ \cad{bba}$.
  \item $\cad{bab}$ ocupa la posición~12.
  \item Antes de las cadenas de longitud $L$ hay $\sum_{k=0}^{L-1}2^{k}=2^{L}-1$ cadenas. Entre las de longitud $L$, el orden lexicográfico con $\cad{a}<\cad{b}$ es el de los números binarios de $L$ cifras, de modo que la cadena de valor $v$ es la $v$-ésima (contando desde~0). Su posición es $2^{L}-1+v$.
\end{enumerate}
\resultado posición de $\cad{bab}$: 12; fórmula general $2^{L}-1+v$.
\comprobacion $\cad{bab}=101_2=5$, $L=3$: $2^{3}-1+5=12$.
\end{solucion}

\begin{solucion}{eje:6-diagonal}
\datos $s_0=10110$, $s_1=00101$, $s_2=11100$, $s_3=01011$, $s_4=11111$ (primeros términos).
\planteamiento leer la diagonal $s_k(k)$ y cambiar cada bit.
\pasos
\begin{enumerate}[label=\arabic*.]
  \item Diagonal: $s_0(0)=1$, $s_1(1)=0$, $s_2(2)=1$, $s_3(3)=1$, $s_4(4)=1$.
  \item $d=0\,1\,0\,0\,0\ldots$
  \item $d$ difiere de $s_k$ en la posición $k$: $d_0=0\neq1$, $d_1=1\neq0$, $d_2=0\neq1$, $d_3=0\neq1$, $d_4=0\neq1$.
\end{enumerate}
\resultado $d=01000\ldots$, distinta de $s_0,\dots,s_4$.
\comprobacion el mismo procedimiento, aplicado a toda la lista, produce una sucesión que no está en ella: es la demostración del \cref{cor:t6-nonum}.
\end{solucion}

\begin{solucion}{eje:6-subconjN}
\datos $S\subseteq\N$ infinito.
\planteamiento definir recursivamente $s_k$ como el $k$-ésimo elemento de $S$; usar que todo subconjunto no vacío de $\N$ tiene mínimo.
\pasos
\begin{enumerate}[label=\arabic*.]
  \item $s_0=\min S$ y $s_{k+1}=\min\bigl(S\setminus\{s_0,\dots,s_k\}\bigr)$; el conjunto es no vacío porque $S$ es infinito.
  \item $s_0<s_1<s_2<\cdots$: $s_{k+1}$ es un elemento de $S$ distinto de los anteriores y mayor que $s_k$ (si fuese menor, habría sido elegido antes). Una sucesión estrictamente creciente es inyectiva.
  \item Sobreyectiva: supóngase $x\in S$ distinto de todos los $s_k$. Entonces, para cada $k$, $x\in S\setminus\{s_0,\dots,s_{k-1}\}$ y por ser $s_k$ el mínimo, $s_k\le x$. Habría infinitos naturales distintos $s_0<s_1<\cdots$ menores o iguales que $x$: absurdo.
\end{enumerate}
\resultado $k\mapsto s_k$ es una biyección $\N\to S$. \hfill$\blacksquare$
\comprobacion con $S=2\N$, $s_k=2k$: es la biyección del \cref{ej:t6-pares}.
\end{solucion}

\begin{solucion}{eje:6-producto}
\datos $g(m,n)=2^{m}(2n+1)-1$.
\planteamiento todo entero positivo se escribe de forma única como potencia de~2 por un impar.
\pasos
\begin{enumerate}[label=\arabic*.]
  \item Sea $N\in\N$. El número $N+1\ge1$ se escribe $N+1=2^{m}\cdot u$ con $u$ impar, dividiendo por~2 mientras sea par; $m$ y $u$ son únicos. Con $u=2n+1$, $n\in\N$, se obtiene $g(m,n)=N$.
  \item La unicidad de $(m,u)$ da la inyectividad, y la existencia, la sobreyectividad: $g$ es biyectiva.
  \item (b) Sean $\alpha\colon\N\to A$ y $\beta\colon\N\to B$ biyecciones. $h(m,n)=(\alpha(m),\beta(n))$ es una biyección $\N\times\N\to A\times B$ (su inversa es $(a,b)\mapsto(\alpha^{-1}(a),\beta^{-1}(b))$). Entonces $h\circ g^{-1}\colon\N\to A\times B$ es biyectiva.
\end{enumerate}
\resultado $g$ es una biyección y el producto de dos conjuntos infinitos numerables es infinito numerable. \hfill$\blacksquare$
\comprobacion $g(1,2)=2\cdot5-1=9$, y $9+1=10=2^{1}\cdot5$ devuelve $(m,n)=(1,2)$.
\end{solucion}

\begin{solucion}{eje:6-Q}
\datos forma reducida $p/q$, única.
\planteamiento inyectar $\Q$ en un conjunto numerable y aplicar la \cref{prop:t6-criterios}.
\pasos
\begin{enumerate}[label=\arabic*.]
  \item $r=p/q\mapsto(p,q)$ es una función $\Q\to\Z\times\N$, bien definida por la unicidad de la forma reducida, e inyectiva porque $(p,q)$ determina $r=p/q$.
  \item $\Z\times\N$ es infinito numerable por el \cref{eje:6-producto}(b), pues $\Z$ y $\N$ lo son.
  \item Componiendo con una biyección $\Z\times\N\to\N$, se obtiene una inyección $\Q\to\N$; por la \cref{prop:t6-criterios}, $\Q$ es numerable.
  \item $\Q\supseteq\N$ es infinito: es infinito numerable.
\end{enumerate}
\resultado $\Q$ es infinito numerable. \hfill$\blacksquare$
\comprobacion $-\tfrac{4}{6}=-\tfrac23\mapsto(-2,3)$, y ningún otro racional va a $(-2,3)$.
\end{solucion}

\begin{solucion}{eje:6-PA}
\datos $A$ conjunto; se supone $h\colon\Pa(A)\to A$ inyectiva.
\planteamiento convertir la inyección en una sobreyección $A\to\Pa(A)$ y aplicar el \cref{teo:t6-cantor}.
\pasos
\begin{enumerate}[label=\arabic*.]
  \item Si $A=\emptyset$, $\Pa(A)=\{\emptyset\}$ no es vacío y no hay ninguna función de un conjunto no vacío en $\emptyset$.
  \item Si $A\neq\emptyset$, $\Pa(A)\neq\emptyset$ y, por la \cref{prop:t4-laterales}, $h$ tiene una inversa por la izquierda $r\colon A\to\Pa(A)$ con $r\circ h=\id_{\Pa(A)}$.
  \item Como $r\circ h$ es sobreyectiva, la función exterior $r$ es sobreyectiva (\cref{prop:t4-reciproco}(ii)).
  \item Eso contradice el \cref{teo:t6-cantor}.
\end{enumerate}
\resultado no existe ninguna inyección $\Pa(A)\to A$: $\abs{\Pa(A)}\not\le\abs{A}$. \hfill$\blacksquare$
\comprobacion para $A$ finito con $n$ elementos se reduce al palomar: $2^{n}>n$.
\end{solucion}

\begin{solucion}{eje:6-intervalos}
\datos $M=\{0\}\cup\{\tfrac1n: n\ge1\}\subseteq[0,1]$; $f(0)=\tfrac12$, $f(\tfrac1n)=\tfrac1{n+2}$, $f(x)=x$ si $x\notin M$.
\planteamiento la idea del «hotel de Hilbert»: desplazar una sucesión dos lugares para hacer sitio a $0$ y a~$1$.
\pasos
\begin{enumerate}[label=\arabic*.]
  \item (a) $2x=2x'\Rightarrow x=x'$; dado $y\in[0,2]$, $x=y/2\in[0,1]$. Es biyectiva.
  \item (b) Imagen de $M$: $f(M)=\{\tfrac12,\tfrac13,\tfrac14,\dots\}=\{\tfrac1k: k\ge2\}$, y $f\colon M\to f(M)$ es biyectiva ($0\mapsto\tfrac12$, $\tfrac1n\mapsto\tfrac1{n+2}$ con $n+2\ge3$).
  \item Imagen del resto: $[0,1]\setminus M$ se envía a sí mismo, y es $(0,1)\setminus\{\tfrac1k:k\ge2\}$, porque $0$ y $1=\tfrac11$ están en $M$.
  \item Las dos imágenes son disjuntas y su unión es $(0,1)$. En cada parte $f$ es inyectiva, luego es inyectiva en total; y cubre $(0,1)$, luego es sobreyectiva.
\end{enumerate}
\resultado $f$ es una biyección $[0,1]\to(0,1)$, y $\abs{[0,1]}=\abs{(0,1)}$. \hfill$\blacksquare$
\comprobacion $f(1)=\tfrac13$, $f(\tfrac12)=\tfrac14$, $f(0{,}7)=0{,}7$; la preimagen de $\tfrac12$ es $0$ y la de $\tfrac15$ es $\tfrac13$.
\end{solucion}
